\chapter{生存分析（二）}
\sourcepages{6}{1-2}
\subsection*{主要内容}
生存分析的基本概念；生存率的估计方法；生存曲线比较的方法；Cox回归分析；生存分析方法的正确应用。
第4章的KM法与log-rank检验只能处理一个分组因素，同时考虑肿瘤分级、大小与是否复发等多个因素时，需采用回归模型。Cox模型仅对风险之比建模，对风险随时间的变化（基线风险$h_0(t)$）不作任何分布假设。估计时，在每个死亡时刻考虑风险集中恰为该个体发生事件的条件概率，$h_0(t)$在分子与分母中相消。以条件化消去冗余参数的方法，与第2章的条件logistic回归相同。

Cox模型的关键假设是比例风险，即任意两个个体的风险之比不随时间变化。本章后半部分的比例风险检验、时依协变量与竞争风险，均针对该假设不成立或问题更为复杂的情形。本章应掌握HR的含义，以及单一HR能够概括效应的条件。
\section{Cox回归分析}
\sourcepages{6}{3-5}
\subsection{问题的提出}
log-rank检验属单因素分析方法，用于不同组间生存曲线的比较，应用条件是影响生存率的各协变量组间均衡可比。如何分析多个因素对生存率（生存曲线）的影响？
\begin{itemize}
\item Logistic回归？不仅考虑结局，还考虑时间。
\item 一般线性模型？生存时间的分布非正态、删失数据。
\end{itemize}
常用的回归模型：指数分布回归模型、Weibull分布回归模型、Cox比例风险回归模型。
\subsection{Cox回归}
\sourcepages{6}{6-7}
Cox比例风险回归模型（Cox's proportional hazards regression model），简称Cox回归。以生存结局和生存时间为因变量，可同时分析众多因素对生存期的影响，分析带有删失生存时间的资料，且不要求资料服从特定的分布类型。英国统计学家D. R. Cox于1972年提出。
\begin{annotation}三大回归：一般线性回归、Logistics回归、Cox回归。\end{annotation}
\subsubsection{Cox回归模型结构}
\[
\underbrace{h(t)}_{t\text{时刻的风险函数}}=
\underbrace{h_0(t)}_{\substack{t\text{时刻的基线风险函数}\\\text{非参数部分}}}
\times\underbrace{\exp(\beta_1X_1+\beta_2X_2+\cdots+\beta_pX_p)}_{\text{风险比，未知参数}}.
\]
半参数模型（semi-parametric model）。
\[
h(t)=h_0(t)\times\exp(\beta_1X_1)\exp(\beta_2X_2)\cdots\exp(\beta_pX_p).
\]
乘法模型。
\begin{annotation}只关心风险比；每个危险因素效应的乘积。\end{annotation}
\subsubsection{模型的假定}
\sourcepages{6}{8-9}
任两个个体在同一时刻的风险函数之比，称为风险比（hazard ratio, HR）。
\begin{align*}
\HR&=\frac{h_i(t)}{h_j(t)}
=\frac{h_0(t)\exp(\beta_1X_{i1}+\beta_2X_{i2}+\cdots+\beta_pX_{ip})}
 {h_0(t)\exp(\beta_1X_{j1}+\beta_2X_{j2}+\cdots+\beta_pX_{jp})}\\
&=\exp[\beta_1(X_{i1}-X_{j1})+\beta_2(X_{i2}-X_{j2})+\cdots+\beta_p(X_{ip}-X_{jp})].
\end{align*}
自变量的效应不随时间而改变，称为比例风险（proportional hazard）假定，简称PH假定。
\[
\ln h_i(t)-\ln h_j(t)=\beta_1(X_{i1}-X_{j1})+\beta_2(X_{i2}-X_{j2})+\cdots+\beta_p(X_{ip}-X_{jp}).
\]
\cfigure{图13\quad 比例风险假定示意图}{\begin{tikzpicture}[x=1cm,y=1cm,font=\small]
\draw (0,0) rectangle (7,3.8);\node[left] at (0,1.9) {$\ln h(t)$};\node[below] at (3.5,0) {time};
\draw[PlotPrimary,line width=1pt] (.7,2.55) .. controls (1.5,3) and (2.2,3.4) .. (2.8,3.2) .. controls (3.5,3) and (4.2,2.6) .. (4.7,2.4);
\draw[PlotSecondary,line width=1pt,dashed] (.7,1.15) .. controls (1.5,1.6) and (2.2,2) .. (2.8,1.8) .. controls (3.5,1.6) and (4.2,1.2) .. (4.7,1);
\node[right] at (5,2.5) {case $i$};\node[right] at (5,1.1) {case $j$};\end{tikzpicture}}
\begin{annotation}$i,j$ 代表两个个体。PH假定下，两人的 $\ln h(t)$ 曲线随时间平行，间距为 $\sum_k\beta_k(X_{ik}-X_{jk})$。

三个效应指标的区别：OR是优势之比（Logistic回归，病例对照、横断面资料）；RR（risk ratio）是固定时间内累积发生概率之比，$\RR(t)=F_i(t)/F_j(t)$（固定随访期的队列、试验）；HR是瞬时风险之比（生存分析）。即使HR恒定，$\RR(t)$ 也随 $t$ 变化，且当事件不罕见时 $\RR(t)$ 比HR更接近1。三者只有在事件罕见时才近似相等，不能混用。\end{annotation}
\subsubsection{参数的解释}
\sourcepages{6}{10}
$\beta_j$：在其它自变量不变条件下，变量 $X_j$ 每增加一个单位所引起的风险比的自然对数。
\[
\ln\HR_j=\beta_j,\qquad\HR_j=\exp(\beta_j).
\]
有的资料同时写作 $\RR_j=\exp(\beta_j)$，这里的“RR”只能理解为风险率之比（rate/hazard ratio），不是固定时点的累积风险之比；本书统一记为HR。
\begin{annotation}解释方式与Logistic回归系数类似（对数尺度上的加法、原尺度上的乘法），但对象不同：Logistic的 $\exp(\beta)$ 是OR，Cox的 $\exp(\beta)$ 是HR，在PH假定下对所有时刻成立。\end{annotation}
\subsubsection{参数的估计和假设检验}
\sourcepages{6}{11-15}
借助部分似然函数（partial likelihood function）用最大似然法估计回归系数。
\begin{annotation}只关心 $\beta$，不关心基线。\end{annotation}
\textbf{部分似然函数的构建}

部分似然函数（Partial Likelihood Function）：不包含基线风险的信息，仅聚焦于事件发生的顺序和协变量的关系。构建部分似然函数的核心思路是“条件风险”：在某一事件发生时间点 $t_k$，考虑所有此时仍处于“风险集”（即尚未发生事件且未删失）的个体，计算实际发生事件的个体恰好是该个体的概率。

对第 $k$ 个事件时间 $t_k$，风险集 $R(t_k)$ 表示在 $t_k$ 前一瞬间仍处于观察中（尚未发生事件、也尚未删失，即 $\widetilde T_l\ge t_k$）的所有个体的集合；恰在 $t_k$ 发生事件或删失的人都包括在内。
\[
P(\text{个体 }i\text{ 在 }t_k\text{ 发生事件}\mid R(t_k)\text{ 中某个个体发生事件})
=\frac{\exp(\beta X_i)}{\sum_{l\in R(t_k)}\exp(\beta X_l)}.
\]
整个样本的部分似然函数就是所有事件发生时间点上该概率的乘积。
\[
L(\beta)=\prod_{k=1}^m\frac{\exp(\beta X_{i_k})}{\sum_{l\in R(t_k)}\exp(\beta X_l)}.
\]
设样本中有 $n$ 个个体，事件发生的时间点按从小到大排序为 $t_1<t_2<t_3<\cdots<t_m$（$m\le n$）。
\[
\ln L(\beta)=\sum_{k=1}^m\left[\beta X_{i_k}-\ln\left(\sum_{l\in R(t_k)}\exp(\beta X_l)\right)\right].
\]
求导并令导数为0，数值求解似然方程组。
\begin{derivation}{部分似然的得分函数与并列事件}
\dpart{假设}PH模型 $h(t\mid\bm x)=h_0(t)\exp(\bm x^\top\bm\beta)$；给定协变量后独立删失；各事件时间互不相同（无并列）。
\dpart{关键等式}在 $t_k$ 时刻，风险集中某人发生事件的条件概率里 $h_0(t_k)$ 在分子分母中相消，部分似然 $L(\bm\beta)=\prod_k e^{\bm x_{i_k}^\top\bm\beta}/\sum_{l\in R(t_k)}e^{\bm x_l^\top\bm\beta}$。对 $\bm\beta$ 求导：
\[
\bm U(\bm\beta)=\frac{\partial\ln L}{\partial\bm\beta}=\sum_{k=1}^m\Bigl\{\bm x_{i_k}-\bar{\bm x}(t_k;\bm\beta)\Bigr\},\qquad
\bar{\bm x}(t_k;\bm\beta)=\frac{\sum_{l\in R(t_k)}\bm x_l\,e^{\bm x_l^\top\bm\beta}}{\sum_{l\in R(t_k)}e^{\bm x_l^\top\bm\beta}} .
\]
$\bar{\bm x}(t_k;\bm\beta)$ 是风险集内以 $e^{\bm x_l^\top\bm\beta}$ 为权的协变量加权均值，即模型下“谁会发生事件”的期望协变量值。信息阵为风险集内的加权协方差阵之和：$\bm I(\bm\beta)=\sum_k\sum_{l\in R(t_k)}w_{kl}\{\bm x_l-\bar{\bm x}(t_k)\}\{\bm x_l-\bar{\bm x}(t_k)\}^\top$，$w_{kl}=e^{\bm x_l^\top\bm\beta}/\sum_{r\in R(t_k)}e^{\bm x_r^\top\bm\beta}$。
\dpart{结论}得分方程 $\bm U(\widehat{\bm\beta})=\bm 0$ 表示：各事件时点上“实际发生事件者的协变量”与“风险集加权均值”之差的总和为0。模型只含一个分组变量时，在 $\bm\beta=\bm 0$ 处的比分检验就是log-rank检验（无并列事件时完全相同）。
\dpart{并列事件}若 $t_k$ 时刻有 $d_k>1$ 人同时发生事件，上式不再精确。常用近似：Breslow法（分母用 $\{\sum_{l\in R}e^{\bm x_l^\top\bm\beta}\}^{d_k}$，SAS PHREG、SPSS默认）、Efron法（逐个扣除已发生者的平均贡献，R \texttt{coxph}默认，更精确）；精确法（exact/discrete）计算量大，适用于并列很多或配对资料（条件Logistic回归即用discrete）。并列少时几种方法结果接近；不同软件的系数可能因此略有差别。
\end{derivation}
假设检验方法类似于Logistic回归，有似然比检验、Wald检验和Score检验，检验统计量在大样本下近似服从卡方分布，自由度为待检验的独立约束个数（参数个数）。

以第4章例3的两组资料拟合仅含肿瘤大小（$\geqslant3$\,cm记为1）的Cox模型，输出如下：
\par\noindent
\begin{coursecode}{R}
library(survival)
d <- data.frame(time=c(t1,t2), status=c(s1,s2), size=rep(0:1,c(18,12)))  # t1,s1,t2,s2见第4章
summary(coxph(Surv(time,status)~size, data=d))
\end{coursecode}
\begin{routput}
  n= 30, number of events= 21 

       coef exp(coef) se(coef)     z Pr(>|z|)  
size 1.2752    3.5793   0.5242 2.433    0.015 *

     exp(coef) exp(-coef) lower .95 upper .95
size     3.579     0.2794     1.281     9.999

Concordance= 0.654  (se = 0.055 )
Likelihood ratio test= 6.18  on 1 df,   p=0.01
Wald test            = 5.92  on 1 df,   p=0.01
Score (logrank) test = 6.71  on 1 df,   p=0.01
\end{routput}
结果解释：肿瘤$\geqslant3$\,cm者的死亡风险为$<3$\,cm者的3.58倍（95\%CI 1.28--10.0）；\texttt{exp(-coef)}为反向比较的HR（0.279）。最后三行为同一假设$H_0:\beta=0$的三种检验，与第2章所述三种检验方法对应。比分检验标注为“logrank”，因为仅含一个分组变量时，比分检验即log-rank检验；此处6.71与\texttt{survdiff}的6.77略有差别，源于24月与29月两处并列死亡的处理方式不同。\texttt{Concordance}为C指数，即随机抽取两人时风险评分较高者先发生死亡的概率，0.5表示无区分能力。

第4章由$(A_1/T_1)/(A_2/T_2)$得到的相对危险度为0.36，其倒数为2.8，与HR 3.58方向一致、大小相近，但仅为粗略近似。
\subsubsection{自变量的筛选策略}
\sourcepages{6}{16}
确定筛选的指标、标准。结合研究目的、专业知识、既往经验：前进法、后退法、逐步回归法；其他：最优组合、结合专业知识的不同模型及其比较。
\subsubsection{Cox回归的主要用途}
\sourcepages{6}{17}
影响因素分析、筛选危险因素；校正协变量后的组间比较；多变量生存预测；生存率估计。
\subsubsection{实例}
\sourcepages{6}{18-23}
\textbf{例5} 30例胃癌患者的随访记录见表8，试进行胃癌肿瘤患者生存情况的影响因素分析。
\ctable{表8\quad 30例胃癌患者生存资料原始记录表}{\footnotesize\begin{tabular}{rrrrrllrrrr}\toprule id&age&grade&size&relapse&start&end&t&status&PI&$S(t)$\\\midrule
1&76&3&0&1&1996/03/27&1996/08/24&5&0&6.174&0.899\\
2&23&2&1&0&1997/12/17&2000/08/03&32&1&5.605&0.102\\
3&52&2&1&0&1998/09/16&2001/02/02&29&1&5.605&0.233\\
4&74&2&1&1&1996/04/29&1997/02/23&10&1&7.040&0.624\\
5$^a$&67&3&0&1&1998/01/16&1998/04/16&3&0&6.174$^a$&0.899$^a$\\
6&47&1&0&0&1997/02/10&2010/12/30&167&0&1.580&0.570\\
7&37&3&0&0&1997/11/14&1998/05/13&6&0&4.739&0.975\\
8$^b$&67&3&0&1&1998/01/16&1998/04/16&3&0&6.174&0.899\\
9&45&3&0&0&1996/05/16&2000/09/22&52&1&4.739&0.276\\
10&73&2&1&1&1996/10/16&1997/05/14&7&0&7.040&0.777\\
$\cdots$&$\cdots$&$\cdots$&$\cdots$&$\cdots$&$\cdots$&$\cdots$&$\cdots$&$\cdots$&$\cdots$&$\cdots$\\
30&29&2&1&0&1996/11/12&1999/05/01&30&1&5.605&0.169\\
\bottomrule\end{tabular}\par\smallskip\begin{minipage}{\linewidth}\footnotesize
$^a$原表此行印作id 8、PI 4.739、$S(t)$ 0.767。按表2，这一记录是5号患者；其协变量（grade 3、size 0、relapse 1）与1号相同，按本表其他各行所用系数，PI应为6.174。原值与7号（PI 4.739、$t=6$、$S=0.975$）矛盾：同一PI下 $S(t)$ 不可能在 $t=3$ 时为0.767、到 $t=6$ 时反而升至0.975。$S(t)$ 按与1号、8号的一致性改为0.899。\\
$^b$8号的记录与5号完全相同，可能是原表抄录重复，原始资料无法核实。\\
PI列由系数约（1.580，2.446，1.435）算得（如 $3\times1.5797=4.739$，$2\times1.5797+2.4456=5.605$，$5.605+1.435=7.040$），与表9的系数（1.572，2.431，1.431）略有出入；按表9系数，1号患者PI为6.147。两组系数可能来自并列事件处理不同的两次拟合。原资料未提供30例的完整资料，表9无法复算。
\end{minipage}}
\ctable{表9\quad 30例胃癌患者多变量Cox回归分析结果}{\begin{tabular}{lrrrrrl}\toprule
Variable&$b$&$S_b$&Wald $\chi^2$&$P$&HR&HR 95\%CI\\\midrule
grade&1.572&0.470&11.193&0.001&4.816&(1.918, 12.097)\\
size&2.431&0.712&11.645&0.001&11.367&(2.814, 45.918)\\
relapse&1.431&0.786&3.311&0.069&4.183&(0.896, 19.537)\\\bottomrule\end{tabular}}
\[
h(t)=h_0(t)\times\exp(1.572\,\mathrm{grade}+2.431\,\mathrm{size}+1.431\,\mathrm{relapse}).
\]
预后指数（prognostic index, PI）：$1.572\,\mathrm{grade}+2.431\,\mathrm{size}+1.431\,\mathrm{relapse}$。
\begin{annotation}
$\HR=\exp(\beta)$。在肿瘤大小、复发情况相同的条件下，肿瘤分级每增加1级，死亡风险（瞬时风险）变为原来的4.816倍（95\%CI 1.918--12.097）。原资料写作“5.357倍”，与表9中 $\exp(1.572)=4.816$ 不符，予以更正。注意grade是以1、2、3赋值的数值变量，这一写法假定II级对I级、III级对II级的HR相同（均为4.816，III级对I级为 $4.816^2=23.2$）；如不愿作此假定，应以哑变量进入模型。size的HR为11.367（95\%CI 2.814--45.918），区间很宽，反映样本只有30例。

预后指数PI越大，死亡风险越高，预后越差，可按PI对患者分层。relapse（是否复发）是随访期间才发生的事件，作为基线变量进入模型有引入未来信息的问题，见“时依协变量”一节。
\end{annotation}
\textbf{$t$ 时刻生存率的估计}
\[
\widehat S(t)=[\widehat S_0(t)]^{\exp(\sum b_jX_j)},
\]
\[
\widehat S_0(t)=\prod_{t_{(i)}\le t}\left[\exp\frac{-d_i}{\sum_{s\in R_i}\exp(\sum b_jX_j)}\right].
\]
$t_{(i)}$ 为排序的生存时间（不含删失值），$d_i$ 为 $t_{(i)}$ 时刻死亡数，$R_i$ 为 $t_{(i)}$ 时刻前一瞬间的风险集，分母中的 $\sum b_jX_j$ 按风险集中每个人 $s$ 各自的协变量计算。
\begin{derivation}{累积基线风险与个体生存率预测}
\dpart{假设}PH模型成立，$h(t\mid\bm x)=h_0(t)\exp(\bm x^\top\bm\beta)$；给定协变量后独立删失。
\dpart{关键等式}对风险函数积分：
\begin{gather*}
H(t\mid\bm x)=\int_0^th_0(u)e^{\bm x^\top\bm\beta}\,du=H_0(t)\,e^{\bm x^\top\bm\beta},\\
S(t\mid\bm x)=e^{-H(t\mid\bm x)}=\{e^{-H_0(t)}\}^{e^{\bm x^\top\bm\beta}}=S_0(t)^{e^{\bm x^\top\bm\beta}} .
\end{gather*}
累积基线风险用Breslow估计：在每个事件时点，把 $d_i$ 例事件按风险集中各人的相对风险 $e^{\bm x_s^\top\widehat{\bm\beta}}$ 分摊，
\[
\widehat H_0(t)=\sum_{t_{(i)}\le t}\frac{d_i}{\sum_{s\in R_i}\exp(\bm x_s^\top\widehat{\bm\beta})},\qquad
\widehat S_0(t)=\exp\{-\widehat H_0(t)\}=\prod_{t_{(i)}\le t}\exp\left\{\frac{-d_i}{\sum_{s\in R_i}\exp(\bm x_s^\top\widehat{\bm\beta})}\right\},
\]
即上面的公式。
\dpart{结论}$\widehat S(t\mid\bm x)=\widehat S_0(t)^{\exp(\bm x^\top\widehat{\bm\beta})}$，其中 $\bm x^\top\widehat{\bm\beta}$ 就是预后指数PI。$\widehat S_0$ 对应“所有协变量取0”的个体；许多软件把协变量中心化，输出的基线生存对应“协变量取均值”的个体，使用时须与PI的计算方式一致。
\dpart{解释}只有 $\widehat{\bm\beta}$ 只能比较相对风险；要预测某人的绝对生存率，还需要估计的基线风险 $\widehat H_0(t)$、该人的协变量值，且只在观察到的随访时间范围内、对与样本相似的人群才可靠。用于临床预测时还应评价区分度和校准度。
\end{derivation}
\section{Cox回归的正确应用}
\sourcepages{6}{24-26}
\subsection{满足比例风险假定（PH假定）}
Cox模型的基本假定是比例风险假定。
\begin{enumerate}
\item 最简单的方法是观察按该变量分组的Kaplan--Meier生存曲线，若生存曲线明显交叉，提示不满足PH假定。
\item 绘制按该变量分组的累积风险函数 $H(t)=\int_0^th(u)\,du$ 对生存时间 $t$ 的关系图。PH假定下 $H_1(t)=H_2(t)\,e^{\beta}$，两条曲线应\emph{成比例}（任一时刻纵坐标之比恒定），而不是平行或等距。比例关系不便目测，故常取对数作图。
\item 专门的假设检验方法。
\end{enumerate}
\textbf{log--log图}\quad 由 $S(t\mid x)=S_0(t)^{e^{\beta x}}$ 得
\[
\ln\{-\ln S(t\mid x)\}=\ln H(t\mid x)=\ln H_0(t)+\beta x .
\]
以 $\ln t$（或 $t$）为横轴、按组绘制 $\ln\{-\ln\widehat S(t)\}$，PH假定成立时各组曲线\emph{平行}，垂直距离约为 $\beta$；曲线交叉或间距随时间明显变化提示违背PH假定。

\textbf{Schoenfeld残差}\quad 对第 $k$ 个事件，$\bm r_k=\bm x_{i_k}-\bar{\bm x}(t_k;\widehat{\bm\beta})$，即发生事件者的协变量减去风险集的加权均值（得分函数中的各项）。若某变量的真实系数随时间变化，$\beta_j(t)$，则标准化Schoenfeld残差对时间作图大致沿 $\beta_j(t)-\widehat\beta_j$ 的形状变化；PH假定下应在0附近水平分布。以残差与时间（或其函数）的相关性作检验，即R的 \texttt{cox.zph}、Stata的 \texttt{estat phtest}，可逐个变量和整体检验。$P>0.05$ 只表示未发现违背PH的证据。
\begin{sikao}[违背比例风险假定的实例]
\texttt{survival}包自带的\texttt{veteran}数据为137例晚期肺癌患者的临床试验资料，协变量包括治疗方式\texttt{trt}、Karnofsky体能评分\texttt{karno}（0--100，越高越好）与年龄：
\par\noindent
\begin{coursecode}{R}
f <- coxph(Surv(time,status)~trt+karno+age, data=veteran)
cox.zph(f)                       # Schoenfeld残差检验
f2 <- coxph(Surv(time,status)~trt+karno+age+tt(karno), data=veteran,
            tt=function(x,t,...) x*log(t))       # karno的效应随log(t)变化
round(summary(f2)$coefficients,4)
exp(coef(f2)["karno"]*10 + coef(f2)["tt(karno)"]*10*log(c(10,100,300)))
\end{coursecode}
\begin{routput}
        chisq df       p
trt     0.284  1 0.59439
karno  12.003  1 0.00053
age     2.100  1 0.14735
GLOBAL 19.102  3 0.00026

             coef exp(coef) se(coef)       z Pr(>|z|)
trt        0.0437    1.0447   0.1924  0.2273   0.8202
karno     -0.0839    0.9195   0.0173 -4.8364   0.0000
age       -0.0040    0.9960   0.0092 -0.4322   0.6656
tt(karno)  0.0133    1.0134   0.0044  3.0500   0.0023
[1] 0.5870032 0.7974291 0.9229409
\end{routput}
\texttt{karno}的比例风险检验$P=0.0005$，比例风险假定不成立。加入$\texttt{karno}\times\log t$项后，交互项有统计学意义（$P=0.002$），表明体能评分的作用随时间减弱：评分高10分对应的HR，第10天为0.59，第100天为0.80，第300天为0.92，作用已接近消失。

原模型给出的单一HR（$e^{-0.0344\times10}=0.71$）是这一变化过程的某种平均，既不代表早期，也不代表后期的效应。这一结果在临床上可以解释：入组时的体能状态对近期死亡有较强的预测作用，随时间推移，病情变化使基线评分的预测作用逐渐减弱。
\end{sikao}
\cfigure{图14\quad 胃癌数据肿瘤大小（size）的PH假定判定（A：生存曲线图；B：累积风险图）}{\begin{tikzpicture}\begin{axis}[width=7.1cm,height=5.1cm,axis lines=left,xmin=0,xmax=185,ymin=0,ymax=103,xlabel={Survival time (months)},ylabel={Survival probability (\%)},title={A: Survival Function},legend style={font=\scriptsize},legend pos=north east]
\addplot[PlotPrimary,line width=1pt,const plot] coordinates {(0,100.000000) (15,92.857143) (20,85.119048) (21,77.380952) (24,69.642857) (29,61.904762) (52,53.061224) (72,44.217687) (94,35.374150) (136,26.530612) (144,17.687075) (177,0.000000)};\addlegendentry{size $<3$ cm}
\addplot[PlotSecondary,line width=1pt,const plot,dashed] coordinates {(0,100.000000) (1.0,91.670000) (2.0,83.330000) (7.0,83.330000) (10.0,74.070000) (12.0,64.810000) (13.0,64.810000) (24.0,54.010000) (29.0,43.210000) (30.0,32.410000) (32.0,21.600000) (61.0,10.800000) (62.0,0.000000)};\addlegendentry{size $\ge3$ cm}
\end{axis}\end{tikzpicture}\hfill\begin{tikzpicture}\begin{axis}[width=7.1cm,height=5.1cm,axis lines=left,xmin=0,xmax=185,ymin=0,ymax=2.4,xlabel={Survival time (months)},ylabel={Cumulative hazard},title={B: Cumulative hazard function}]
\addplot[PlotPrimary,line width=1pt,const plot] coordinates {(0,0) (15,0.075) (20,0.15) (21,0.24) (24,0.34) (29,0.48) (52,0.64) (72,0.82) (94,1.03) (136,1.32) (144,1.72) (177,2.3)};
\addplot[PlotSecondary,line width=1pt,const plot,dashed] coordinates {(0,0) (1,0.085) (2,0.18) (10,0.29) (12,0.43) (24,0.64) (29,0.86) (30,1.14) (32,1.54) (61,2.25) (62,2.33)};
\end{axis}\end{tikzpicture}}
\subsection{时依协变量}
\sourcepages{6}{27-29}
当资料不满足比例风险（PH）假定时：分段拟合Cox回归模型；或扩展模型，让协变量或系数随时间变化。应区分两种扩展：
\begin{itemize}
\item \textbf{时依协变量} $X(t)$：协变量的\emph{取值}随时间变化，系数不变，$h(t)=h_0(t)\exp\{\beta X(t)\}$。$t$ 时刻的风险只取决于当时（或之前）的 $X$ 值。
\item \textbf{时变系数} $\beta(t)$：协变量取值固定，其\emph{效应}随时间变化，$h(t)=h_0(t)\exp\{\beta(t)X\}$，常用 $\beta(t)=\beta_0+\beta_1g(t)$ 形式，即在模型中加入 $X\times g(t)$ 交互项。它是检验和处理违背PH假定的直接方法。
\end{itemize}
时依协变量又分为外在（external，又称外生）和内在（internal，又称内生）两类：外在时依协变量的变化路径不受个体自身事件过程的影响，与个体是否存活无关，如年龄、日历时间、环境污染水平、预先安排的治疗剂量；内在时依协变量是在个体身上测得、只有个体存活并在随访中才有定义的指标，如血压、肿瘤复发、实验室检查结果，它们本身可能处于暴露与结局的因果通路上，解释时要谨慎。有的资料把“取值不随时间改变、HR随时间改变”称为外在时依协变量，实际描述的是时变系数。

\textbf{例6} 研究1945年后接触核辐射的妇女患乳腺癌的风险。
\[
h(t)=h_0(t)e^{1.05x_E-0.05x_Et}.
\]
$x_E=1,0$（暴露与否）是\emph{固定}的协变量，HR $=\exp(1.05-0.05t)$ 随时间下降，属于时变系数模型（$x_E$ 与 $t$ 的交互）：$t=0$ 时HR为2.86，$t=1,2,3$ 年时依次为2.72、2.59、2.46、……。
\begin{annotation}只有在PH假定成立时，才能用一个HR概括暴露效应；本例的HR随时间变化，应报告HR随 $t$ 的变化或分时段的HR。\end{annotation}
\textbf{例7} 随访开始时是吸烟的（$X_E(t)=1$，当 $t<t_0$ 时），随访到某时刻 $t_0$ 开始戒烟（$X_E(t)=0$，当 $t\ge t_0$ 时）。吸烟状态是时依协变量（个体行为，属内在协变量），系数 $\beta$ 不随时间变化，而同一个人在 $t_0$ 前后的风险不同；两个吸烟史不同的人的风险之比随时间变化，在这个意义上不再满足“比例风险”。
\begin{annotation}实现方法是把每个人的随访拆成若干段 $(\text{start},\text{stop}]$，每段内协变量取值固定。\end{annotation}
\textbf{避免引入未来信息}\quad 表8中的relapse（是否复发）是手术后随访期间才发生的事件，表1只记录了“是否复发”，没有复发日期，无法核实其记录时点。若把“随访中曾复发”当作手术时就已存在的基线协变量，复发者在复发前的那段时间也被当作“已复发”，而能够复发的人必须先活到复发，这段“必然存活”的时间会使复发组显得风险偏低（不朽时间偏倚），HR有偏。正确的做法是记录复发时间，把relapse作为时依协变量。例如4号患者（$t=10$ 月死亡）若在第6个月复发，应拆为两行：
\begin{center}\small\begin{tabular}{rrrrrr}\toprule id&start&stop&relapse&status&说明\\\midrule
4&0&6&0&0&复发前，未发生死亡\\
4&6&10&1&1&复发后，第10个月死亡\\\bottomrule\end{tabular}\end{center}
R中用 \texttt{coxph(Surv(start, stop, status) \textasciitilde{} grade + size + relapse)} 拟合，每个时刻只使用当时已知的复发状态。“第6个月复发”只是示例，原始资料需补充复发时间后才能重新分析。
\subsection{样本量估算}
\sourcepages{6}{30-32}
\begin{enumerate}
\item 经验估算方法：阳性结局事件数至少为模型参数个数的10--20倍。这是经验规则而非由检验效能推出的结果；参数按实际自由度计数（$c$ 类分类变量计 $c-1$ 个，样条、交互项另计），而不是按原始变量个数计。
\item 考察某感兴趣变量 $X_1$ 主效应的情况，即检验 $\beta_1$ 是否等于0。
\end{enumerate}
Cox回归的检验效能由\emph{事件数}决定，而不是由总例数决定。先算所需事件数
\[
D=\frac{(Z_{\alpha/2}+Z_\beta)^2}{(1-R^2)\sigma^2(\ln\HR)^2},
\]
再按研究结束时的预期事件发生率 $P$ 换算为总样本量
\[
n=\frac DP=\frac{(Z_{\alpha/2}+Z_\beta)^2}{P(1-R^2)\sigma^2(\ln\HR)^2}.
\]
$P$ 为研究结束时终点事件发生率；$R^2$ 为主效应变量 $X_1$ 对其他协变量作回归分析时的决定系数（$1/(1-R^2)$ 即方差膨胀因子）；$\sigma^2$ 为 $X_1$ 的方差；$\ln\HR$ 为 $X_1$ 每增加1个单位的对数风险比，即回归系数 $\beta_1$。$\sigma^2$ 与 $\ln\HR$ 必须按同一编码计算。

\textbf{例8} 拟进行一项临床试验对一新药与标准药的抗肺癌疗效进行比较。假定研究结束时死亡率 $P=70\%$，服用药物的标准差为0.8，对数风险比为0.5（即标准药死亡风险是新药的 $e^{0.5}=1.65$ 倍），服用药物对其他协变量作回归分析时的决定系数估计值为0.30，试问按照检验水准0.05、检验功效0.90，需要观察多少例肺癌患者？

原题解法：
\[
n=\frac{(1.96+1.282)^2}{0.70\times(1-0.30)\times0.8^2\times0.5^2}\approx135.
\]
\revnote{“服用药物”是0/1二分类变量（标准药0、新药1），若服新药的比例为 $p$，则 $\Var(X)=p(1-p)\le1/4$，标准差最大为0.5，不可能为0.8，故上式不成立。按 $1:1$ 分配，$p=0.5$，$\sigma^2=0.25$：
\[
D=\frac{(1.96+1.282)^2}{(1-0.30)\times0.25\times0.5^2}=240.2\approx241\ \text{例死亡},\qquad
n=\frac{240.2}{0.70}=343.1\approx344\ \text{例}.
\]
若按 $2:1$ 分配（$p=2/3$），$\sigma^2=2/9$，所需例数增至约386。此外，随机分组时处理与基线协变量独立，$R^2$ 应接近0；取 $R^2=0$ 时 $D\approx169$，$n\approx241$。题设 $R^2=0.30$ 更符合观察性研究中暴露与协变量相关的情形。}
\subsection{其他模型}
\sourcepages{6}{33-37}
\subsubsection{竞争风险模型（competing risk model）}
结局包括多种状态，不仅仅是生存和死亡，或者死亡原因有多种。例如，$Y=1$（癌症）、2（心血管疾病）、0（删失）。分析危险因素对患者死于不同原因的影响大小。Fine \& Gray模型（1999年）。

\textbf{例子} 某研究人员收集了2000年确诊为轻度认知损害（MCI）的518例老年患者资料，包括基本人口学特征、生活方式、体格检查和合并疾病信息等，并于2025年签完成25次随访调查，主要观察结局为是否发生阿尔兹海默病（AD）。问题：影响MCI向AD转归的因素有哪些？

如果MCI患者在观察期间死于癌症、心血管疾病等原因而未发生AD，就不能为AD的发病做出贡献，即死亡“竞争”了AD的发生。传统统计方法将发生AD前死亡的个体、失访个体和未发生AD个体均按删失数据（censored data）处理：用来估计AD的原因别风险是可以的，但用 $1-\mathrm{KM}$ 估计“实际有多少人会发生AD”会偏高（见下）。
\cfigure{一般生存过程与竞争风险过程}{\begin{tikzpicture}[every node/.style={draw,minimum width=16mm,font=\small},node distance=7mm and 9mm]
\node (a) {0};\node[right=of a,yshift=6mm] (e) {事件1};\node[right=of a,yshift=-6mm] (s) {右删失};\draw[PlotPrimary,line width=.65pt,-{Stealth}] (a)--(e);\draw[PlotPrimary,line width=.65pt,-{Stealth}] (a)--(s);\node[draw=none,below=of a,xshift=15mm] {一般生存过程};
\node[right=32mm of e,yshift=-6mm] (b) {0};\node[right=of b,yshift=12mm] (e1) {事件1};\node[right=of b] (e2) {事件2};\node[right=of b,yshift=-17mm] (c) {右删失};\node[draw=none] at ($(e2)!.5!(c)$) {$\vdots$};\draw[PlotPrimary,line width=.65pt,-{Stealth}] (b)--(e1);\draw[PlotPrimary,line width=.65pt,-{Stealth}] (b)--(e2);\draw[PlotPrimary,line width=.65pt,-{Stealth}] (b)--(c);\node[draw=none,below=21mm of b,xshift=15mm] {竞争风险过程};\end{tikzpicture}}
\cfigure{图2\quad 两种方法估计的复发累积发生率比较（$1-\mathrm{KM}$ 与CIF）}{\begin{tikzpicture}\begin{axis}[width=11cm,height=6cm,axis lines=left,xmin=12,xmax=120,ymin=0,ymax=.9,xtick={12,24,...,120},ytick={0,.1,...,.9},xlabel={生存时间/月},ylabel={累积风险率},legend pos=north west]
\addplot[PlotPrimary,line width=1pt] coordinates {(12,.10)(24,.25)(36,.34)(48,.41)(60,.45)(72,.54)(84,.54)(96,.59)(108,.59)(120,.80)};\addlegendentry{生存分析}
\addplot[PlotSecondary,dashed,line width=1pt] coordinates {(12,.10)(24,.17)(36,.21)(48,.23)(60,.30)(72,.35)(84,.35)(96,.40)(108,.40)(120,.57)};\addlegendentry{竞争风险}
\end{axis}\end{tikzpicture}}
删失只是阻止了你观察到目标结局事件，而竞争事件则是阻止了目标结局事件的发生，所以删失并不是竞争事件。在竞争风险模型中，一个关键点是，一个受试者只能发生一个结局事件类型，这样该事件的发生就排除了后续再发生其他类型事件。
\begin{derivation}{累积发生函数（CIF）与两类回归模型}
\dpart{记号}$T$ 为首个事件的时间，$J\in\{1,\ldots,K\}$ 为事件类型。原因别风险 $h_k(t)=\lim_{\Delta t\to0}P(t\le T<t+\Delta t,J=k\mid T\ge t)/\Delta t$；总风险 $h(t)=\sum_kh_k(t)$，无任何事件的生存函数 $S(t)=\exp\{-\int_0^th(u)\,du\}$。
\dpart{关键等式}第 $k$ 类事件的累积发生函数
\[
F_k(t)=P(T\le t,\ J=k)=\int_0^tS(u^-)\,h_k(u)\,du,\qquad \sum_kF_k(t)=1-S(t).
\]
在 $u$ 时刻发生第 $k$ 类事件，要求此前未发生\emph{任何}事件（$S(u^-)$），且在 $u$ 时刻发生的是第 $k$ 类（$h_k(u)$）。
\dpart{结论}(1) 把其他类事件当删失，KM法估计的是 $S_k^*(t)=\exp\{-\int_0^th_k(u)\,du\}$，$1-S_k^*(t)$ 只在“竞争事件不存在”的假想世界中才是发生概率。由于 $S(u)\le S_k^*(u)$，$1-S_k^*(t)\ge F_k(t)$，即 $1-\mathrm{KM}$ 高估实际的累积发生率，竞争事件越多高估越严重（上图“生存分析”曲线高于“竞争风险”曲线即此）。$F_k$ 应用Aalen--Johansen估计（R \texttt{cmprsk::cuminc}、\texttt{survfit} 的多状态形式）。(2) 回归模型有两类：
\end{derivation}
\ctable{原因别Cox模型与Fine--Gray模型}{\small\begin{tabularx}{\textwidth}{p{2cm}XX}\toprule
&原因别Cox模型（cause-specific hazard）&Fine--Gray模型（subdistribution hazard）\\\midrule
建模对象&原因别风险 $h_k(t\mid\bm x)$&子分布风险 $\lambda_k(t)=-\dfrac{d}{dt}\ln\{1-F_k(t)\}$\\\addlinespace
风险集&尚未发生任何事件且仍在观察的人；发生竞争事件者退出（按删失处理）&除尚未发生事件者外，已发生竞争事件者\emph{仍保留}在风险集中（权重按删失分布调整）\\\addlinespace
$\exp(\beta)$ 的含义&原因别HR：在仍未发生任何事件的人中，协变量对第 $k$ 类事件瞬时发生率的作用&子分布HR：协变量对第 $k$ 类事件累积发生率 $F_k(t)$ 的作用方向和大小\\\addlinespace
适用问题&病因研究（暴露是否直接影响该事件的发生速率）&预测与卫生决策（某类人最终有多少会发生该事件）\\\bottomrule\end{tabularx}}
同一变量的两类HR可以方向不同：例如某因素提高死亡风险，即使它不影响AD的原因别风险，也会因“死得早”而降低AD的累积发生率。报告时宜同时给出两类模型，或说明选择依据。若研究关心的是全因死亡，则各种死因都是同一终点事件，不存在竞争风险；只有区分死因、关心特定死因时才需要竞争风险方法。
\subsubsection{限制性立方样条函数}
\sourcepages{6}{38-39}
限制性立方样条函数（Restricted Cubic Spline, RCS）。分析暴露因素（连续型定量变量）和结局之间是否存在非线性关系、剂量反应关系。分段多项式，强制要求立方样条函数在连续性自变量 $x$ 的两端区间内成线性函数。
\Needspace{5\baselineskip}
\par\noindent
\begin{coursecode}{R}
library(rms)
dd <- datadist(data); options(datadist="dd")
fit <- cph(Surv(time,death) ~ rcs(age,4) + sex, data=data)
anova(fit)                         # 总体关联与非线性的检验
plot(Predict(fit, age, ref.zero=TRUE, fun=exp))   # HR(age)，参照值默认为中位数
\end{coursecode}
\par
\begin{derivation}{RCS的基函数、自由度与HR曲线}
\dpart{构造}取 $k$ 个结点 $t_1<\cdots<t_k$（常按 $x$ 的分位数放置，$k=4$ 时取5\%、35\%、65\%、95\%分位数）。记 $(u)_+=\max(u,0)$，RCS为
\[
f(x)=\beta_1x+\sum_{j=1}^{k-2}\beta_{j+1}S_j(x),\quad
S_j(x)=(x-t_j)_+^3-(x-t_{k-1})_+^3\frac{t_k-t_j}{t_k-t_{k-1}}+(x-t_k)_+^3\frac{t_{k-1}-t_j}{t_k-t_{k-1}}.
\]
各段是三次多项式，在结点处函数值、一阶、二阶导数连续；在 $t_1$ 以左和 $t_k$ 以右，三次项和二次项相互抵消，$f$ 为直线（“限制性”即指两端线性），避免尾部数据少时曲线剧烈摆动。
\dpart{自由度}$k$ 个结点的RCS在Cox模型中占 $k-1$ 个参数（$x$ 的线性项1个、非线性基函数 $k-2$ 个；Cox模型无截距）。代码中 \texttt{rcs(age,4)} 即4个结点、3个参数。
\dpart{两种检验}(1) 总体关联：$H_0:\beta_1=\cdots=\beta_{k-1}=0$，$\mathrm{df}=k-1$，检验 $x$ 与风险是否有关；(2) 非线性：$H_0:\beta_2=\cdots=\beta_{k-1}=0$，$\mathrm{df}=k-2$，检验关系是否偏离线性。非线性检验 $P>0.05$ 只说明没有偏离线性的充分证据，此时可考虑用线性项；总体关联不显著时，曲线的形状不宜解释。
\dpart{HR曲线}以 $x_{\mathrm{ref}}$ 为参照，
\[
\HR(x)=\exp\{f(x)-f(x_{\mathrm{ref}})\}=\exp(\bm c^\top\widehat{\bm\beta}),\qquad \bm c=\bm b(x)-\bm b(x_{\mathrm{ref}}),
\]
$\bm b(x)=(x,S_1(x),\ldots,S_{k-2}(x))^\top$ 为基函数向量。由系数协方差阵 $\widehat{\bm V}$ 得 $\Var\{\ln\widehat{\HR}(x)\}=\bm c^\top\widehat{\bm V}\bm c$，95\%CI为 $\exp\{\bm c^\top\widehat{\bm\beta}\pm1.96\sqrt{\bm c^\top\widehat{\bm V}\bm c}\}$。
\dpart{解释}(1) 参照值决定曲线的上下位置：$x=x_{\mathrm{ref}}$ 处HR恒为1、置信区间宽度为0，越远离参照值区间越宽；改变参照值会使整条曲线上下平移，但形状不变。参照值应事先说明（如中位数、临床正常值，图中竖线50岁即参照值的位置）。(2) 两端数据稀少，置信区间很宽，尾部的“上翘”“下降”可能只是随机波动，应结合数据分布（如叠加直方图）解释。(3) 结点个数宜事先确定（3--5个），不应反复试换以求显著。
\end{derivation}
\cfigure{Death Risk}{\begin{tikzpicture}\begin{axis}[width=11cm,height=7cm,axis lines=left,xmin=17,xmax=80,ymin=0,ymax=4.45,xlabel={Age},ylabel={HR (95\%CI)},ytick={1,2,3,4}]
\addplot[draw=none,fill=PlotPrimary!14] coordinates {(19.731801,1.840467) (19.961686,1.813230) (20.191571,1.778210) (20.421456,1.750973) (20.651341,1.723735) (20.881226,1.696498) (21.111111,1.669261) (21.340996,1.642023) (21.570881,1.614786) (21.800766,1.587549) (22.030651,1.560311) (22.260536,1.536965) (22.490421,1.509728) (22.720307,1.486381) (22.950192,1.459144) (23.180077,1.435798) (23.409962,1.412451) (23.639847,1.392996) (23.869732,1.369650) (24.099617,1.346304) (24.329502,1.322957) (24.559387,1.303502) (24.789272,1.284047) (25.019157,1.264591) (25.249042,1.245136) (25.478927,1.221790) (25.708812,1.202335) (25.938697,1.182879) (26.168582,1.163424) (26.398467,1.147860) (26.628352,1.124514) (26.858238,1.108949) (27.088123,1.093385) (27.318008,1.073930) (27.547893,1.058366) (27.777778,1.038911) (28.007663,1.027237) (28.237548,1.011673) (28.467433,0.996109) (28.697318,0.980545) (28.927203,0.964981) (29.157088,0.949416) (29.386973,0.933852) (29.616858,0.918288) (29.846743,0.906615) (30.076628,0.891051) (30.306513,0.879377) (30.536398,0.863813) (30.766284,0.852140) (30.996169,0.840467) (31.226054,0.828794) (31.455939,0.817121) (31.685824,0.805447) (31.915709,0.793774) (32.145594,0.782101) (32.375479,0.774319) (32.605364,0.766537) (32.835249,0.754864) (33.065134,0.747082) (33.295019,0.735409) (33.524904,0.727626) (33.754789,0.719844) (33.984674,0.712062) (34.214559,0.704280) (34.444444,0.696498) (34.674330,0.692607) (34.904215,0.684825) (35.134100,0.680934) (35.363985,0.673152) (35.593870,0.669261) (35.823755,0.665370) (36.053640,0.661479) (36.283525,0.657588) (36.513410,0.653696) (36.743295,0.649805) (36.973180,0.649805) (37.203065,0.645914) (37.432950,0.645914) (37.662835,0.642023) (37.892720,0.642023) (38.122605,0.642023) (38.352490,0.642023) (38.582375,0.642023) (38.812261,0.642023) (39.042146,0.642023) (39.272031,0.642023) (39.501916,0.642023) (39.731801,0.645914) (39.961686,0.645914) (40.191571,0.645914) (40.421456,0.649805) (40.651341,0.653696) (40.881226,0.657588) (41.111111,0.661479) (41.340996,0.661479) (41.570881,0.665370) (41.800766,0.673152) (42.030651,0.677043) (42.260536,0.680934) (42.490421,0.684825) (42.720307,0.692607) (42.950192,0.696498) (43.180077,0.700389) (43.409962,0.708171) (43.639847,0.712062) (43.869732,0.719844) (44.099617,0.727626) (44.329502,0.735409) (44.559387,0.743191) (44.789272,0.750973) (45.019157,0.758755) (45.249042,0.766537) (45.478927,0.774319) (45.708812,0.785992) (45.938697,0.793774) (46.168582,0.801556) (46.398467,0.813230) (46.628352,0.824903) (46.858238,0.836576) (47.088123,0.844358) (47.318008,0.856031) (47.547893,0.867704) (47.777778,0.879377) (48.007663,0.891051) (48.237548,0.902724) (48.467433,0.914397) (48.697318,0.926070) (48.927203,0.937743) (49.157088,0.953307) (49.386973,0.964981) (49.616858,0.976654) (49.846743,0.996109) (50.076628,1.019455) (50.306513,1.042802) (50.536398,1.073930) (50.766284,1.108949) (50.996169,1.143969) (51.226054,1.175097) (51.455939,1.210117) (51.685824,1.245136) (51.915709,1.284047) (52.145594,1.319066) (52.375479,1.361868) (52.605364,1.396887) (52.835249,1.439689) (53.065134,1.478599) (53.295019,1.513619) (53.524904,1.556420) (53.754789,1.595331) (53.984674,1.634241) (54.214559,1.673152) (54.444444,1.712062) (54.674330,1.750973) (54.904215,1.785992) (55.134100,1.824903) (55.363985,1.859922) (55.593870,1.898833) (55.823755,1.933852) (56.053640,1.972763) (56.283525,2.003891) (56.513410,2.038911) (56.743295,2.073930) (56.973180,2.108949) (57.203065,2.143969) (57.432950,2.175097) (57.662835,2.210117) (57.892720,2.241245) (58.122605,2.272374) (58.352490,2.303502) (58.582375,2.330739) (58.812261,2.357977) (59.042146,2.385214) (59.272031,2.416342) (59.501916,2.439689) (59.731801,2.466926) (59.961686,2.494163) (60.191571,2.521401) (60.421456,2.548638) (60.651341,2.568093) (60.881226,2.591440) (61.111111,2.614786) (61.340996,2.638132) (61.570881,2.661479) (61.800766,2.680934) (62.030651,2.704280) (62.260536,2.723735) (62.490421,2.739300) (62.720307,2.762646) (62.950192,2.782101) (63.180077,2.801556) (63.409962,2.817121) (63.639847,2.836576) (63.869732,2.856031) (64.099617,2.875486) (64.329502,2.891051) (64.559387,2.906615) (64.789272,2.926070) (65.019157,2.945525) (65.249042,2.961089) (65.478927,2.980545) (65.708812,2.996109) (65.938697,3.015564) (66.168582,3.031128) (66.398467,3.046693) (66.628352,3.066148) (66.858238,3.085603) (67.088123,3.105058) (67.318008,3.124514) (67.547893,3.140078) (67.777778,3.159533) (68.007663,3.178988) (68.237548,3.198444) (68.467433,3.217899) (68.697318,3.237354) (68.927203,3.256809) (69.157088,3.276265) (69.386973,3.299611) (69.616858,3.322957) (69.846743,3.342412) (70.076628,3.365759) (70.306513,3.389105) (70.536398,3.412451) (70.766284,3.435798) (70.996169,3.459144) (71.226054,3.486381) (71.455939,3.509728) (71.685824,3.536965) (71.915709,3.560311) (72.145594,3.587549) (72.375479,3.614786) (72.605364,3.638132) (72.835249,3.669261) (73.065134,3.692607) (73.295019,3.723735) (73.524904,3.750973) (73.754789,3.782101) (73.984674,3.805447) (74.214559,3.840467) (74.444444,3.871595) (74.674330,3.898833) (74.904215,3.933852) (75.134100,3.961089) (75.363985,3.996109) (75.593870,4.027237) (75.823755,4.058366) (76.053640,4.093385) (76.283525,4.124514) (76.513410,4.159533) (76.743295,4.194553) (76.973180,4.229572) (76.973180,1.704280) (76.743295,1.708171) (76.513410,1.715953) (76.283525,1.719844) (76.053640,1.719844) (75.823755,1.723735) (75.593870,1.727626) (75.363985,1.735409) (75.134100,1.739300) (74.904215,1.739300) (74.674330,1.743191) (74.444444,1.747082) (74.214559,1.750973) (73.984674,1.754864) (73.754789,1.754864) (73.524904,1.762646) (73.295019,1.766537) (73.065134,1.770428) (72.835249,1.770428) (72.605364,1.770428) (72.375479,1.774319) (72.145594,1.778210) (71.915709,1.782101) (71.685824,1.782101) (71.455939,1.785992) (71.226054,1.789883) (70.996169,1.789883) (70.766284,1.793774) (70.536398,1.793774) (70.306513,1.797665) (70.076628,1.797665) (69.846743,1.801556) (69.616858,1.801556) (69.386973,1.801556) (69.157088,1.805447) (68.927203,1.805447) (68.697318,1.805447) (68.467433,1.805447) (68.237548,1.805447) (68.007663,1.805447) (67.777778,1.805447) (67.547893,1.801556) (67.318008,1.801556) (67.088123,1.801556) (66.858238,1.797665) (66.628352,1.797665) (66.398467,1.793774) (66.168582,1.789883) (65.938697,1.789883) (65.708812,1.785992) (65.478927,1.782101) (65.249042,1.774319) (65.019157,1.770428) (64.789272,1.766537) (64.559387,1.762646) (64.329502,1.758755) (64.099617,1.750973) (63.869732,1.743191) (63.639847,1.739300) (63.409962,1.727626) (63.180077,1.723735) (62.950192,1.715953) (62.720307,1.704280) (62.490421,1.700389) (62.260536,1.692607) (62.030651,1.680934) (61.800766,1.673152) (61.570881,1.665370) (61.340996,1.653696) (61.111111,1.645914) (60.881226,1.634241) (60.651341,1.622568) (60.421456,1.614786) (60.191571,1.599222) (59.961686,1.591440) (59.731801,1.579767) (59.501916,1.568093) (59.272031,1.556420) (59.042146,1.540856) (58.812261,1.529183) (58.582375,1.517510) (58.352490,1.505837) (58.122605,1.494163) (57.892720,1.482490) (57.662835,1.466926) (57.432950,1.455253) (57.203065,1.443580) (56.973180,1.428016) (56.743295,1.416342) (56.513410,1.400778) (56.283525,1.385214) (56.053640,1.373541) (55.823755,1.361868) (55.593870,1.346304) (55.363985,1.330739) (55.134100,1.319066) (54.904215,1.303502) (54.674330,1.287938) (54.444444,1.276265) (54.214559,1.260700) (53.984674,1.245136) (53.754789,1.229572) (53.524904,1.217899) (53.295019,1.202335) (53.065134,1.186770) (52.835249,1.171206) (52.605364,1.159533) (52.375479,1.143969) (52.145594,1.128405) (51.915709,1.116732) (51.685824,1.101167) (51.455939,1.085603) (51.226054,1.073930) (50.996169,1.058366) (50.766284,1.046693) (50.536398,1.035019) (50.306513,1.023346) (50.076628,1.011673) (49.846743,0.976654) (49.616858,0.953307) (49.386973,0.922179) (49.157088,0.898833) (48.927203,0.867704) (48.697318,0.844358) (48.467433,0.817121) (48.237548,0.793774) (48.007663,0.770428) (47.777778,0.747082) (47.547893,0.723735) (47.318008,0.704280) (47.088123,0.688716) (46.858238,0.665370) (46.628352,0.645914) (46.398467,0.630350) (46.168582,0.614786) (45.938697,0.599222) (45.708812,0.583658) (45.478927,0.568093) (45.249042,0.552529) (45.019157,0.540856) (44.789272,0.529183) (44.559387,0.517510) (44.329502,0.505837) (44.099617,0.498054) (43.869732,0.486381) (43.639847,0.478599) (43.409962,0.470817) (43.180077,0.463035) (42.950192,0.459144) (42.720307,0.451362) (42.490421,0.439689) (42.260536,0.435798) (42.030651,0.431907) (41.800766,0.428016) (41.570881,0.420233) (41.340996,0.412451) (41.111111,0.412451) (40.881226,0.408560) (40.651341,0.404669) (40.421456,0.400778) (40.191571,0.392996) (39.961686,0.392996) (39.731801,0.389105) (39.501916,0.385214) (39.272031,0.385214) (39.042146,0.385214) (38.812261,0.381323) (38.582375,0.377432) (38.352490,0.373541) (38.122605,0.373541) (37.892720,0.373541) (37.662835,0.369650) (37.432950,0.365759) (37.203065,0.365759) (36.973180,0.361868) (36.743295,0.361868) (36.513410,0.357977) (36.283525,0.357977) (36.053640,0.354086) (35.823755,0.350195) (35.593870,0.350195) (35.363985,0.350195) (35.134100,0.346304) (34.904215,0.346304) (34.674330,0.342412) (34.444444,0.342412) (34.214559,0.338521) (33.984674,0.338521) (33.754789,0.334630) (33.524904,0.330739) (33.295019,0.330739) (33.065134,0.330739) (32.835249,0.326848) (32.605364,0.322957) (32.375479,0.322957) (32.145594,0.322957) (31.915709,0.319066) (31.685824,0.315175) (31.455939,0.315175) (31.226054,0.311284) (30.996169,0.311284) (30.766284,0.307393) (30.536398,0.307393) (30.306513,0.303502) (30.076628,0.299611) (29.846743,0.299611) (29.616858,0.299611) (29.386973,0.295720) (29.157088,0.291829) (28.927203,0.287938) (28.697318,0.287938) (28.467433,0.287938) (28.237548,0.284047) (28.007663,0.280156) (27.777778,0.280156) (27.547893,0.280156) (27.318008,0.276265) (27.088123,0.272374) (26.858238,0.268482) (26.628352,0.268482) (26.398467,0.268482) (26.168582,0.264591) (25.938697,0.260700) (25.708812,0.260700) (25.478927,0.260700) (25.249042,0.256809) (25.019157,0.252918) (24.789272,0.249027) (24.559387,0.249027) (24.329502,0.249027) (24.099617,0.245136) (23.869732,0.241245) (23.639847,0.241245) (23.409962,0.237354) (23.180077,0.237354) (22.950192,0.233463) (22.720307,0.233463) (22.490421,0.229572) (22.260536,0.225681) (22.030651,0.225681) (21.800766,0.225681) (21.570881,0.221790) (21.340996,0.221790) (21.111111,0.217899) (20.881226,0.217899) (20.651341,0.214008) (20.421456,0.214008) (20.191571,0.210117) (19.961686,0.206226) (19.731801,0.206226)}\closedcycle;
\addplot[PlotPrimary,line width=1pt] coordinates {(19.731801,0.624514) (19.961686,0.624514) (20.191571,0.620623) (20.421456,0.616732) (20.651341,0.616732) (20.881226,0.614786) (21.111111,0.612840) (21.340996,0.608949) (21.570881,0.605058) (21.800766,0.601167) (22.030651,0.599222) (22.260536,0.597276) (22.490421,0.597276) (22.720307,0.593385) (22.950192,0.589494) (23.180077,0.587549) (23.409962,0.585603) (23.639847,0.585603) (23.869732,0.581712) (24.099617,0.577821) (24.329502,0.573930) (24.559387,0.573930) (24.789272,0.573930) (25.019157,0.570039) (25.249042,0.566148) (25.478927,0.564202) (25.708812,0.562257) (25.938697,0.562257) (26.168582,0.558366) (26.398467,0.554475) (26.628352,0.552529) (26.858238,0.550584) (27.088123,0.550584) (27.318008,0.546693) (27.547893,0.542802) (27.777778,0.540856) (28.007663,0.538911) (28.237548,0.538911) (28.467433,0.535019) (28.697318,0.531128) (28.927203,0.531128) (29.157088,0.531128) (29.386973,0.527237) (29.616858,0.523346) (29.846743,0.523346) (30.076628,0.523346) (30.306513,0.519455) (30.536398,0.515564) (30.766284,0.515564) (30.996169,0.511673) (31.226054,0.511673) (31.455939,0.507782) (31.685824,0.507782) (31.915709,0.503891) (32.145594,0.501946) (32.375479,0.500000) (32.605364,0.500000) (32.835249,0.496109) (33.065134,0.496109) (33.295019,0.496109) (33.524904,0.496109) (33.754789,0.492218) (33.984674,0.492218) (34.214559,0.488327) (34.444444,0.488327) (34.674330,0.488327) (34.904215,0.488327) (35.134100,0.488327) (35.363985,0.484436) (35.593870,0.484436) (35.823755,0.484436) (36.053640,0.484436) (36.283525,0.484436) (36.513410,0.484436) (36.743295,0.484436) (36.973180,0.484436) (37.203065,0.484436) (37.432950,0.488327) (37.662835,0.488327) (37.892720,0.488327) (38.122605,0.488327) (38.352490,0.492218) (38.582375,0.492218) (38.812261,0.492218) (39.042146,0.494163) (39.272031,0.496109) (39.501916,0.500000) (39.731801,0.500000) (39.961686,0.503891) (40.191571,0.503891) (40.421456,0.507782) (40.651341,0.511673) (40.881226,0.515564) (41.111111,0.523346) (41.340996,0.523346) (41.570881,0.531128) (41.800766,0.535019) (42.030651,0.538911) (42.260536,0.546693) (42.490421,0.550584) (42.720307,0.558366) (42.950192,0.564202) (43.180077,0.570039) (43.409962,0.577821) (43.639847,0.585603) (43.869732,0.593385) (44.099617,0.599222) (44.329502,0.608949) (44.559387,0.618677) (44.789272,0.632296) (45.019157,0.642023) (45.249042,0.651751) (45.478927,0.665370) (45.708812,0.677043) (45.938697,0.692607) (46.168582,0.704280) (46.398467,0.715953) (46.628352,0.731518) (46.858238,0.743191) (47.088123,0.760700) (47.318008,0.778210) (47.547893,0.793774) (47.777778,0.809339) (48.007663,0.828794) (48.237548,0.842412) (48.467433,0.861868) (48.697318,0.881323) (48.927203,0.902724) (49.157088,0.922179) (49.386973,0.943580) (49.616858,0.966926) (49.846743,0.986381) (50.076628,1.015564) (50.306513,1.031128) (50.536398,1.052529) (50.766284,1.075875) (50.996169,1.103113) (51.226054,1.124514) (51.455939,1.145914) (51.685824,1.169261) (51.915709,1.196498) (52.145594,1.223735) (52.375479,1.249027) (52.605364,1.274319) (52.835249,1.297665) (53.065134,1.324903) (53.295019,1.350195) (53.524904,1.377432) (53.754789,1.402724) (53.984674,1.426070) (54.214559,1.453307) (54.444444,1.476654) (54.674330,1.501946) (54.904215,1.527237) (55.134100,1.550584) (55.363985,1.573930) (55.593870,1.597276) (55.823755,1.620623) (56.053640,1.643969) (56.283525,1.667315) (56.513410,1.690661) (56.743295,1.714008) (56.973180,1.737354) (57.203065,1.758755) (57.432950,1.780156) (57.662835,1.799611) (57.892720,1.819066) (58.122605,1.842412) (58.352490,1.863813) (58.582375,1.883268) (58.812261,1.898833) (59.042146,1.918288) (59.272031,1.937743) (59.501916,1.955253) (59.731801,1.974708) (59.961686,1.992218) (60.191571,2.007782) (60.421456,2.027237) (60.651341,2.040856) (60.881226,2.058366) (61.111111,2.073930) (61.340996,2.087549) (61.570881,2.105058) (61.800766,2.116732) (62.030651,2.132296) (62.260536,2.147860) (62.490421,2.159533) (62.720307,2.169261) (62.950192,2.180934) (63.180077,2.194553) (63.409962,2.208171) (63.639847,2.221790) (63.869732,2.233463) (64.099617,2.241245) (64.329502,2.252918) (64.559387,2.264591) (64.789272,2.274319) (65.019157,2.284047) (65.249042,2.291829) (65.478927,2.301556) (65.708812,2.311284) (65.938697,2.321012) (66.168582,2.328794) (66.398467,2.338521) (66.628352,2.346304) (66.858238,2.356031) (67.088123,2.361868) (67.318008,2.373541) (67.547893,2.379377) (67.777778,2.385214) (68.007663,2.394942) (68.237548,2.400778) (68.467433,2.412451) (68.697318,2.418288) (68.927203,2.424125) (69.157088,2.429961) (69.386973,2.439689) (69.616858,2.445525) (69.846743,2.453307) (70.076628,2.459144) (70.306513,2.468872) (70.536398,2.474708) (70.766284,2.480545) (70.996169,2.490272) (71.226054,2.496109) (71.455939,2.503891) (71.685824,2.509728) (71.915709,2.519455) (72.145594,2.527237) (72.375479,2.533074) (72.605364,2.542802) (72.835249,2.548638) (73.065134,2.554475) (73.295019,2.564202) (73.524904,2.570039) (73.754789,2.577821) (73.984674,2.585603) (74.214559,2.593385) (74.444444,2.601167) (74.674330,2.608949) (74.904215,2.616732) (75.134100,2.624514) (75.363985,2.632296) (75.593870,2.640078) (75.823755,2.647860) (76.053640,2.655642) (76.283525,2.663424) (76.513410,2.671206) (76.743295,2.678988) (76.973180,2.686770)};
\draw[dashed] (axis cs:17,1)--(axis cs:80,1);\draw[thin] (axis cs:50,0)--(axis cs:50,4.4);\end{axis}\end{tikzpicture}}
\subsection{建模步骤与策略}
\sourcepages{6}{40}
\begin{enumerate}
\item 模型的选择：Cox模型、其他模型（指数模型、Weibull模型）。
\item 建模的步骤：参数估计、假设检验、模型评估、模型筛选。
\item 从单因素分析开始——部分多因素——多因素。
\item 深入思考：变量类型、条件的检验、变量间的关系、拟合优度。
\item 结果解释：统计学意义与专业意义、设计的好坏。
\end{enumerate}
第2章“按研究目的确定建模方案”中的原则同样适用于Cox回归，另需注意：(1) 有效样本量是事件数，每个参数约需10--20个事件；(2) 连续变量的函数形式在 $\ln h$ 尺度上检查（鞅残差图、样条）；(3) 对每个主要变量检查PH假定（log--log图、Schoenfeld残差），违背时采用分层、时变系数或分时段模型；(4) 协变量必须在时间起点已知，随访中才发生的事件（如复发）应作时依协变量；(5) 预测模型用C统计量评价区分度，用校准曲线评价校准度，并作内部验证。
\section{案例分析}
\sourcepages{6}{41-45}
\subsection{案例介绍与研究设计}
队列研究中，探讨危险因素（risk factor）对死亡风险（mortality risk）的影响：长期业余时间体育锻炼强度和全死亡及死因别死亡的关联，一项前瞻性队列研究。

\textit{Long-Term Leisure-Time Physical Activity Intensity and All-Cause and Cause-Specific Mortality: A Prospective Cohort of US Adults}。

Dong Hoon Lee, Leandro F M Rezende, Hee-Kyung Joh, NaNa Keum, Gerson Ferrari, Juan Pablo Rey-Lopez, Eric B Rimm, Fred K Tabung, Edward L Giovannucci. Circulation. 2022 Aug 16;146(7):523--534. doi:10.1161/CIRCULATIONAHA.121.058162. Epub 2022 Jul 25. PMID:35876019. PMCID:PMC9378548。

题目解析：暴露？结局？对照？设计方法？体育锻炼强度；全死亡/死因别死亡；对照（内对照/外对照/总人口对照/多重对照）；前瞻性队列研究。

\textbf{背景与目的} 目前指南建议的运动水平为每周150到300分钟的中等强度运动（MPA），或每周75至150分钟的剧烈运动（VPA），或等量的中等强度和高强度有氧活动组合。问题：每周运动多长时间，死亡风险最低？
\subsection{数据特点}
\sourcepages{6}{46-47}
本研究数据来自两项大型前瞻性队列研究，一个为护士健康研究（NHS），另一个为医疗专业人员随访研究（HPFS）。NHS研究始于1976年，有121701名女性（年龄30--55岁）登记参与了研究；HPFS研究始于1986年，有51529名男性（年龄40--76岁）登记参与了研究。参与者每两年完成一次关于人口统计学、生活方式和病史信息的问卷调查。

根据代谢当量（MET）将运动分为中等强度运动（$<6$ MET）和剧烈运动（$\ge6$ MET）。中等强度运动包括散步、举重和健美操；剧烈运动包括慢跑、游泳、骑自行车和其他有氧运动。

参与者每2年完成一份随访研究问卷，随访率均超过90\%。本研究对以上两个队列的数据进行了分析，随访时间为1988--2018年，对116221名成年人的运动数据进行分析。
\begin{center}\begin{tikzpicture}[x=.5cm,y=.5cm,line width=1.7pt,line cap=round,color=PlotPrimary]
\fill (0,3.5) circle (.25);\draw (0,3.1)--(-.2,1.9)--(-.8,.2);\draw (-.2,1.9)--(.8,.2);\draw (0,2.9)--(-.8,2.3)--(-1.5,2.7);\draw (0,2.9)--(.7,2.3)--(1.2,2.5);
\begin{scope}[xshift=4cm]\fill (0,3.5) circle (.25);\draw (0,3.1)--(0,1.8)--(-.9,.2);\draw (0,1.8)--(1,1);\draw (0,2.9)--(-1,2.5);\draw (0,2.9)--(.9,2.2)--(1.5,2.7);\draw (1.6,3.1) ellipse (.3 and .45);\end{scope}
\begin{scope}[xshift=8cm]\fill (0,1.7) circle (.25);\draw (-1.3,.9)--(-.3,1.2)--(.8,1.1);\draw (-.2,1.2)--(.7,2.4);\draw (-1.6,.3) .. controls (-1,.8) and (-.5,-.1) .. (0,.3) .. controls (.4,.8) and (.9,-.1) .. (1.5,.3);\end{scope}
\end{tikzpicture}\end{center}
\subsection{方法应用}
\sourcepages{6}{48-49}
\textbf{Cox回归模型的建立}
\begin{itemize}
\item 人年的计算：从基线1988开始，到死亡或随访结束。结果：在2984545人年的随访中（中位随访26年），共有47596名参与者死亡。
\item 使用Cox比例风险回归模型估计体育运动强度与全因死亡率以及特定原因死亡率（CVD死亡和非CVD死亡）之间的关联。
\item 计算关联强度：风险比（HR）和95\%CI。
\item 多因素Cox回归模型：协变量调整种族、CVD家族史、肿瘤家族史、绝经后激素使用、饮酒、能量摄入、吸烟、睡眠、膳食指数、BMI。
\end{itemize}
Cox回归中，采用限制性立方样条函数（Restricted cubic spline, RCS），拟合连续性自变量和事件风险之间的剂量反应关系、非线性关系。
\begin{annotation}增加一个非线性项，可以研究危险因素和结局有没有线性关系/剂量反应关系。\end{annotation}
\subsection{结果及解释}
\sourcepages{6}{50-52}
\ctable{Table 3\quad Associations Between Long-Term Leisure-Time VPA and All-Cause and Cause-Specific Mortality (Pooled Results of HPFS and NHS, 1988--2018)}{\footnotesize\begin{tabular}{lrrll}\toprule Outcome&Deaths&Person-years&Age-adjusted&Multivariable-adjusted\\ &&&HR (95\%CI)&HR (95\%CI)$^*$\\\midrule
\multicolumn{5}{l}{All-cause mortality}\\\multicolumn{5}{l}{VPA, min/wk}\\
0&1194&91 573&1 (Reference)&1 (Reference)\\
1--74&31 297&1 776 378&0.63 (0.59--0.67)&0.87 (0.82--0.93)\\
75--149&7907&539 846&0.50 (0.47--0.53)&0.81 (0.76--0.87)\\
150--224&3494&269 767&0.48 (0.45--0.52)&0.79 (0.74--0.85)\\
225--299&1595&122 439&0.46 (0.43--0.50)&0.77 (0.72--0.84)\\
300--374&937&79 642&0.47 (0.43--0.51)&0.78 (0.72--0.85)\\
375--449&459&39 875&0.43 (0.38--0.48)&0.73 (0.65--0.81)\\
450--599&444&38 387&0.45 (0.40--0.50)&0.76 (0.68--0.85)\\
$\ge600$&269&26 638&0.45 (0.39--0.51)&0.74 (0.65--0.85)\\
\midrule
\multicolumn{5}{l}{CVD mortality}\\\multicolumn{5}{l}{VPA, min/wk}\\
0&285&92 503&1 (Reference)&1 (Reference)\\
1--74&6896&1 805 615&0.56 (0.50--0.64)&0.76 (0.67--0.86)\\
75--149&1625&547 572&0.44 (0.39--0.50)&0.69 (0.60--0.78)\\
150--224&788&273 018&0.45 (0.40--0.52)&0.73 (0.64--0.85)\\
225--299&341&123 927&0.40 (0.34--0.47)&0.67 (0.57--0.79)\\
300--374&203&80 485&0.39 (0.33--0.47)&0.68 (0.56--0.82)\\
375--449&105&40 288&0.38 (0.31--0.48)&0.65 (0.52--0.82)\\
450--599&99&38 787&0.37 (0.29--0.47)&0.67 (0.53--0.84)\\
\midrule
\end{tabular}}
与没有进行剧烈运动的参与者相比，符合指南推荐强度的参与者（即75--149分钟/周），全因死亡率降低19\%（HR=0.81；95\%CI 0.76--0.87），心血管疾病死亡率降低31\%（HR=0.69；95\%CI 0.60--0.78）。
\cfigure{Fig. Joint association of long-term leisure-time VPA and MPA with all-cause and cause-specific mortality}{\begin{tikzpicture}\begin{axis}[width=14.8cm,height=7cm,xmin=-.5,xmax=4.5,ymin=0,ymax=2,ytick={0,.2,...,2},xtick={0,1,2,3,4},xticklabels={0--19,20--149,150--299,300--599,$\ge600$},xlabel={MPA (min/wk)},ylabel={HR (95\%CI)},title={All-cause mortality},axis lines=left,legend style={font=\scriptsize},legend pos=north east]
\draw[thin] (axis cs:-.5,1)--(axis cs:4.5,1);
\addplot[only marks,mark=*,mark size=1.9pt,PlotNeutral,error bars/.cd,y dir=both,y explicit] coordinates {(-0.15000,1.00943) += (0,0.00000) -= (0,0.00000) (0.85000,0.91981) += (0,0.13679) -= (0,0.10849) (1.85000,0.75943) += (0,0.15094) -= (0,0.12736) (2.85000,0.55660) += (0,0.15094) -= (0,0.11321) (3.85000,0.72170) += (0,0.53774) -= (0,0.26415) };\addlegendentry{No VPA}
\addplot[only marks,mark=square*,mark size=1.9pt,PlotSecondary,error bars/.cd,y dir=both,y explicit] coordinates {(-0.07500,0.88679) += (0,0.08962) -= (0,0.09434) (0.92500,0.76415) += (0,0.08019) -= (0,0.08019) (1.92500,0.69340) += (0,0.07075) -= (0,0.05660) (2.92500,0.61792) += (0,0.08019) -= (0,0.06604) (3.92500,0.58491) += (0,0.07075) -= (0,0.07547) };\addlegendentry{1--74}
\addplot[only marks,mark=triangle*,mark size=1.9pt,PlotThird,error bars/.cd,y dir=both,y explicit] coordinates {(0.00000,0.78302) += (0,0.12264) -= (0,0.12264) (1.00000,0.70283) += (0,0.07547) -= (0,0.08491) (2.00000,0.65566) += (0,0.06132) -= (0,0.08019) (3.00000,0.58491) += (0,0.07075) -= (0,0.07075) (4.00000,0.56132) += (0,0.08491) -= (0,0.08019) };\addlegendentry{75--149}
\addplot[only marks,mark=diamond*,mark size=1.9pt,PlotFourth,error bars/.cd,y dir=both,y explicit] coordinates {(0.07500,0.62736) += (0,0.10849) -= (0,0.10377) (1.07500,0.67453) += (0,0.06604) -= (0,0.07547) (2.07500,0.63208) += (0,0.07547) -= (0,0.07547) (3.07500,0.59906) += (0,0.06604) -= (0,0.07075) (4.07500,0.53302) += (0,0.08491) -= (0,0.08491) };\addlegendentry{150--299}
\addplot[only marks,mark=pentagon*,mark size=1.9pt,PlotPrimary,error bars/.cd,y dir=both,y explicit] coordinates {(0.15000,0.57075) += (0,0.14151) -= (0,0.12264) (1.15000,0.58491) += (0,0.07547) -= (0,0.07547) (2.15000,0.62736) += (0,0.08491) -= (0,0.08491) (3.15000,0.60377) += (0,0.08019) -= (0,0.07075) (4.15000,0.58491) += (0,0.12264) -= (0,0.09906) };\addlegendentry{$\ge300$}
\end{axis}\end{tikzpicture}}
剧烈运动 $\ge300$ 分钟/周的参与者，死亡率没有再进一步降低。更长时间的MPA（$\ge600$ 分钟/周）与死亡率的相关性，与300--599分钟/周与死亡率的相关性结果类似。
\cfigure{Figure 1. Dose-response relationship of long-term leisure-time MPA and VPA with all-cause mortality (pooled results of HPFS and NHS, 1988--2018)}{\begin{tikzpicture}\begin{axis}[width=7cm,height=5.4cm,xmin=0,xmax=900,ymin=.4,ymax=1.2,xtick={0,150,300,450,600,750,900},ytick={.4,.6,.8,1,1.2},xlabel={Physical Activity (min/week)},ylabel={HR (95\%CI)},title={All-cause mortality},axis lines=left,title style={font=\small\bfseries},legend style={font=\scriptsize},legend pos=south west]
\addplot[PlotSecondary,line width=1pt,dashed] coordinates {(0.000000,1.000000) (19.285714,0.978133) (25.714286,0.967467) (32.142857,0.958933) (38.571429,0.949333) (45.000000,0.940800) (51.428571,0.932267) (57.857143,0.923733) (64.285714,0.916267) (70.714286,0.909867) (77.142857,0.902400) (83.571429,0.894933) (90.000000,0.889600) (96.428571,0.883200) (102.857143,0.878933) (109.285714,0.867200) (115.714286,0.856533) (122.142857,0.852267) (128.571429,0.846933) (135.000000,0.844800) (141.428571,0.840533) (147.857143,0.836267) (154.285714,0.829867) (160.714286,0.825600) (167.142857,0.823467) (173.571429,0.818133) (180.000000,0.816000) (186.428571,0.809600) (192.857143,0.808533) (199.285714,0.805333) (205.714286,0.801067) (212.142857,0.796800) (218.571429,0.795733) (225.000000,0.793600) (231.428571,0.789333) (237.857143,0.787200) (244.285714,0.785067) (250.714286,0.782933) (257.142857,0.779733) (263.571429,0.777600) (270.000000,0.776533) (276.428571,0.774400) (282.857143,0.772267) (289.285714,0.770133) (295.714286,0.769067) (302.142857,0.768000) (308.571429,0.765867) (315.000000,0.763733) (321.428571,0.762667) (327.857143,0.759467) (334.285714,0.758400) (340.714286,0.758400) (347.142857,0.757333) (353.571429,0.755200) (360.000000,0.754133) (366.428571,0.752000) (372.857143,0.750933) (379.285714,0.747733) (385.714286,0.748800) (392.142857,0.747733) (398.571429,0.745600) (405.000000,0.742400) (411.428571,0.741333) (417.857143,0.742400) (424.285714,0.740267) (430.714286,0.740267) (437.142857,0.738133) (443.571429,0.737067) (450.000000,0.736000) (456.428571,0.734933) (462.857143,0.733867) (469.285714,0.731733) (475.714286,0.729600) (482.142857,0.728533) (488.571429,0.728533) (495.000000,0.727467) (501.428571,0.725333) (507.857143,0.724267) (514.285714,0.723200) (520.714286,0.722133) (527.142857,0.721067) (533.571429,0.720000) (540.000000,0.717867) (546.428571,0.716800) (552.857143,0.715733) (559.285714,0.714667) (565.714286,0.711467) (572.142857,0.709333) (578.571429,0.708267) (585.000000,0.709333) (591.428571,0.708267) (597.857143,0.707200) (604.285714,0.706133) (610.714286,0.705067) (617.142857,0.701867) (623.571429,0.702933) (630.000000,0.700800) (636.428571,0.699733) (642.857143,0.697600) (649.285714,0.697600) (655.714286,0.695467) (662.142857,0.694400) (668.571429,0.693333) (675.000000,0.692267) (681.428571,0.690133) (687.857143,0.689067) (694.285714,0.689067) (700.714286,0.688000) (707.142857,0.685867) (713.571429,0.685867) (720.000000,0.682667) (726.428571,0.681600) (732.857143,0.681600) (739.285714,0.680533) (745.714286,0.679467) (752.142857,0.677333) (758.571429,0.677333) (765.000000,0.676267) (771.428571,0.673067) (777.857143,0.673067) (784.285714,0.672000) (790.714286,0.670933) (797.142857,0.669867) (803.571429,0.668800) (810.000000,0.667733) (816.428571,0.666667) (822.857143,0.665600) (829.285714,0.664533) (835.714286,0.662400) (842.142857,0.662400) (848.571429,0.658133) (855.000000,0.658133) (861.428571,0.658133) (867.857143,0.657067) (874.285714,0.656000) (880.714286,0.654933) (887.142857,0.652800) (893.571429,0.651733) (900.000000,0.651733)};\addlegendentry{Moderate PA, $P$ for non-linear $<0.001$}
\addplot[PlotPrimary,line width=1pt] coordinates {(0.000000,1.000000) (19.285714,0.966400) (25.714286,0.955733) (32.142857,0.945067) (38.571429,0.935467) (45.000000,0.926933) (51.428571,0.918400) (57.857143,0.910933) (64.285714,0.903467) (70.714286,0.897067) (77.142857,0.890667) (83.571429,0.885333) (90.000000,0.878933) (96.428571,0.874667) (102.857143,0.870400) (109.285714,0.867200) (115.714286,0.862933) (122.142857,0.858667) (128.571429,0.856533) (135.000000,0.854400) (141.428571,0.852267) (147.857143,0.850133) (154.285714,0.849067) (160.714286,0.846933) (167.142857,0.845867) (173.571429,0.844800) (180.000000,0.843733) (186.428571,0.842667) (192.857143,0.842667) (199.285714,0.841600) (205.714286,0.841600) (212.142857,0.841600) (218.571429,0.840533) (225.000000,0.840533) (231.428571,0.840533) (237.857143,0.839467) (244.285714,0.839467) (250.714286,0.839467) (257.142857,0.838400) (263.571429,0.838400) (270.000000,0.837333) (276.428571,0.837333) (282.857143,0.837333) (289.285714,0.836267) (295.714286,0.836267) (302.142857,0.836267) (308.571429,0.835200) (315.000000,0.835200) (321.428571,0.835200) (327.857143,0.834133) (334.285714,0.834133) (340.714286,0.834133) (347.142857,0.833067) (353.571429,0.833067) (360.000000,0.833067) (366.428571,0.833067) (372.857143,0.833067) (379.285714,0.832000) (385.714286,0.832000) (392.142857,0.832000) (398.571429,0.830933) (405.000000,0.828800) (411.428571,0.828800) (417.857143,0.828800) (424.285714,0.828800) (430.714286,0.828800) (437.142857,0.827733) (443.571429,0.827733) (450.000000,0.827733) (456.428571,0.826667) (462.857143,0.826667) (469.285714,0.826667) (475.714286,0.825600) (482.142857,0.825600) (488.571429,0.826667) (495.000000,0.824533) (501.428571,0.824533) (507.857143,0.824533) (514.285714,0.824533) (520.714286,0.824533) (527.142857,0.823467) (533.571429,0.823467) (540.000000,0.823467) (546.428571,0.822400) (552.857143,0.822400) (559.285714,0.822400) (565.714286,0.820267) (572.142857,0.820267) (578.571429,0.820267) (585.000000,0.820267) (591.428571,0.820267) (597.857143,0.820267) (604.285714,0.819200) (610.714286,0.819200) (617.142857,0.819200) (623.571429,0.818133) (630.000000,0.818133) (636.428571,0.818133) (642.857143,0.817067) (649.285714,0.817067) (655.714286,0.817067) (662.142857,0.816000) (668.571429,0.816000) (675.000000,0.816000) (681.428571,0.816000) (687.857143,0.816000) (694.285714,0.813867) (700.714286,0.813867) (707.142857,0.813867) (713.571429,0.813867) (720.000000,0.812800) (726.428571,0.812800) (732.857143,0.811733) (739.285714,0.811733) (745.714286,0.811733) (752.142857,0.811733) (758.571429,0.811733) (765.000000,0.811733) (771.428571,0.810667) (777.857143,0.810667) (784.285714,0.810667) (790.714286,0.809600) (797.142857,0.809600) (803.571429,0.809600) (810.000000,0.808533) (816.428571,0.808533) (822.857143,0.808533) (829.285714,0.807467) (835.714286,0.807467) (842.142857,0.807467) (848.571429,0.807467) (855.000000,0.807467) (861.428571,0.805333) (867.857143,0.805333) (874.285714,0.805333) (880.714286,0.805333) (887.142857,0.803200) (893.571429,0.803200) (900.000000,0.804267)};\addlegendentry{Vigorous PA, $P$ for non-linear $<0.001$}
\end{axis}\end{tikzpicture}\hfill\begin{tikzpicture}\begin{axis}[width=7cm,height=5.4cm,xmin=0,xmax=900,ymin=.4,ymax=1.2,xtick={0,150,300,450,600,750,900},ytick={.4,.6,.8,1,1.2},xlabel={Physical Activity (min/week)},ylabel={HR (95\%CI)},title={CVD mortality},axis lines=left,title style={font=\small\bfseries},legend style={font=\scriptsize},legend pos=south west]
\addplot[PlotSecondary,line width=1pt,dashed] coordinates {(0.000000,1.000000) (19.251337,0.974933) (25.668449,0.965333) (32.085561,0.954667) (38.502674,0.944000) (44.919786,0.934400) (51.336898,0.925867) (57.754011,0.916267) (64.171123,0.907733) (70.588235,0.900267) (77.005348,0.892800) (83.422460,0.885333) (89.839572,0.878933) (96.256684,0.866133) (102.673797,0.869333) (109.090909,0.855467) (115.508021,0.844800) (121.925134,0.840533) (128.342246,0.836267) (134.759358,0.829867) (141.176471,0.825600) (147.593583,0.821333) (154.010695,0.816000) (160.427807,0.809600) (166.844920,0.805333) (173.262032,0.802133) (179.679144,0.797867) (186.096257,0.793600) (192.513369,0.789333) (198.930481,0.785067) (205.347594,0.782933) (211.764706,0.778667) (218.181818,0.773333) (224.598930,0.771200) (231.016043,0.769067) (237.433155,0.765867) (243.850267,0.763733) (250.267380,0.760533) (256.684492,0.758400) (263.101604,0.755200) (269.518717,0.752000) (275.935829,0.749867) (282.352941,0.748800) (288.770053,0.745600) (295.187166,0.743467) (301.604278,0.741333) (308.021390,0.739200) (314.438503,0.739200) (320.855615,0.737067) (327.272727,0.734933) (333.689840,0.732800) (340.106952,0.732800) (346.524064,0.730667) (352.941176,0.729600) (359.358289,0.727467) (365.775401,0.726400) (372.192513,0.725333) (378.609626,0.723200) (385.026738,0.721067) (391.443850,0.720000) (397.860963,0.718933) (404.278075,0.716800) (410.695187,0.713600) (417.112299,0.712533) (423.529412,0.710400) (429.946524,0.710400) (436.363636,0.709333) (442.780749,0.708267) (449.197861,0.706133) (455.614973,0.704000) (462.032086,0.701867) (468.449198,0.701867) (474.866310,0.700800) (481.283422,0.698667) (487.700535,0.697600) (494.117647,0.696533) (500.534759,0.693333) (506.951872,0.691200) (513.368984,0.691200) (519.786096,0.689067) (526.203209,0.688000) (532.620321,0.686933) (539.037433,0.683733) (545.454545,0.684800) (551.871658,0.682667) (558.288770,0.680533) (564.705882,0.679467) (571.122995,0.678400) (577.540107,0.677333) (583.957219,0.675200) (590.374332,0.674133) (596.791444,0.673067) (603.208556,0.670933) (609.625668,0.669867) (616.042781,0.667733) (622.459893,0.665600) (628.877005,0.664533) (635.294118,0.663467) (641.711230,0.662400) (648.128342,0.659200) (654.545455,0.658133) (660.962567,0.658133) (667.379679,0.656000) (673.796791,0.656000) (680.213904,0.653867) (686.631016,0.652800) (693.048128,0.651733) (699.465241,0.650667) (705.882353,0.647467) (712.299465,0.647467) (718.716578,0.646400) (725.133690,0.644267) (731.550802,0.643200) (737.967914,0.641067) (744.385027,0.641067) (750.802139,0.638933) (757.219251,0.637867) (763.636364,0.635733) (770.053476,0.633600) (776.470588,0.631467) (782.887701,0.632533) (789.304813,0.631467) (795.721925,0.629333) (802.139037,0.628267) (808.556150,0.627200) (814.973262,0.626133) (821.390374,0.624000) (827.807487,0.621867) (834.224599,0.621867) (840.641711,0.620800) (847.058824,0.618667) (853.475936,0.616533) (859.893048,0.616533) (866.310160,0.615467) (872.727273,0.613333) (879.144385,0.612267) (885.561497,0.611200) (891.978610,0.610133) (898.395722,0.609067)};\addlegendentry{Moderate PA, $P$ for non-linear $<0.001$}
\addplot[PlotPrimary,line width=1pt] coordinates {(0.000000,1.000000) (19.251337,0.963200) (25.668449,0.951467) (32.085561,0.940800) (38.502674,0.930133) (44.919786,0.919467) (51.336898,0.909867) (57.754011,0.902400) (64.171123,0.894933) (70.588235,0.887467) (77.005348,0.881067) (83.422460,0.874667) (89.839572,0.870400) (96.256684,0.866133) (102.673797,0.860800) (109.090909,0.856533) (115.508021,0.854400) (121.925134,0.850133) (128.342246,0.848000) (134.759358,0.845867) (141.176471,0.843733) (147.593583,0.843733) (154.010695,0.841600) (160.427807,0.841600) (166.844920,0.840533) (173.262032,0.840533) (179.679144,0.840533) (186.096257,0.840533) (192.513369,0.840533) (198.930481,0.840533) (205.347594,0.840533) (211.764706,0.840533) (218.181818,0.841600) (224.598930,0.841600) (231.016043,0.842667) (237.433155,0.842667) (243.850267,0.842667) (250.267380,0.842667) (256.684492,0.843733) (263.101604,0.843733) (269.518717,0.843733) (275.935829,0.844800) (282.352941,0.844800) (288.770053,0.845867) (295.187166,0.845867) (301.604278,0.846933) (308.021390,0.846933) (314.438503,0.848000) (320.855615,0.849067) (327.272727,0.849067) (333.689840,0.850133) (340.106952,0.850133) (346.524064,0.850133) (352.941176,0.850133) (359.358289,0.851200) (365.775401,0.851200) (372.192513,0.852267) (378.609626,0.852267) (385.026738,0.853333) (391.443850,0.853333) (397.860963,0.854400) (404.278075,0.854400) (410.695187,0.854400) (417.112299,0.854400) (423.529412,0.854400) (429.946524,0.856533) (436.363636,0.856533) (442.780749,0.857600) (449.197861,0.857600) (455.614973,0.858667) (462.032086,0.858667) (468.449198,0.859733) (474.866310,0.859733) (481.283422,0.859733) (487.700535,0.859733) (494.117647,0.860800) (500.534759,0.860800) (506.951872,0.860800) (513.368984,0.861867) (519.786096,0.862933) (526.203209,0.864000) (532.620321,0.862933) (539.037433,0.864000) (545.454545,0.865067) (551.871658,0.865067) (558.288770,0.865067) (564.705882,0.866133) (571.122995,0.866133) (577.540107,0.867200) (583.957219,0.867200) (590.374332,0.868267) (596.791444,0.868267) (603.208556,0.869333) (609.625668,0.868267) (616.042781,0.869333) (622.459893,0.869333) (628.877005,0.870400) (635.294118,0.870400) (641.711230,0.871467) (648.128342,0.872533) (654.545455,0.873600) (660.962567,0.873600) (667.379679,0.873600) (673.796791,0.874667) (680.213904,0.874667) (686.631016,0.875733) (693.048128,0.875733) (699.465241,0.875733) (705.882353,0.876800) (712.299465,0.877867) (718.716578,0.877867) (725.133690,0.877867) (731.550802,0.877867) (737.967914,0.878933) (744.385027,0.878933) (750.802139,0.880000) (757.219251,0.880000) (763.636364,0.880000) (770.053476,0.882133) (776.470588,0.882133) (782.887701,0.882133) (789.304813,0.882133) (795.721925,0.883200) (802.139037,0.883200) (808.556150,0.884267) (814.973262,0.884267) (821.390374,0.885333) (827.807487,0.885333) (834.224599,0.886400) (840.641711,0.886400) (847.058824,0.886400) (853.475936,0.886400) (859.893048,0.887467) (866.310160,0.887467) (872.727273,0.889600) (879.144385,0.889600) (885.561497,0.890667) (891.978610,0.890667) (898.395722,0.890667)};\addlegendentry{Vigorous PA, $P$ for non-linear $<0.001$}
\end{axis}\end{tikzpicture}\par\medskip\begin{tikzpicture}\begin{axis}[width=7cm,height=5.4cm,xmin=0,xmax=900,ymin=0,ymax=28,xtick={0,150,300,450,600,750,900},ytick={0,5,10,15,20,25},xlabel={Physical Activity (min/week)},ylabel={Percent},title={Distribution of moderate and vigorous PA},title style={font=\scriptsize},axis lines=left]
\fill[PlotPrimary!75] (axis cs:0.7660,0.0706) rectangle (axis cs:24.7660,27.0480);
\fill[PlotPrimary!75] (axis cs:26.2979,0.0706) rectangle (axis cs:50.2979,18.6441);
\fill[PlotPrimary!75] (axis cs:51.8298,0.0706) rectangle (axis cs:75.8298,12.7119);
\fill[PlotSecondary!35] (axis cs:51.8298,12.7825) rectangle (axis cs:75.8298,13.1356);
\fill[PlotPrimary!75] (axis cs:77.3617,0.0706) rectangle (axis cs:101.3617,9.1808);
\fill[PlotSecondary!35] (axis cs:77.3617,9.2514) rectangle (axis cs:101.3617,9.3220);
\fill[PlotPrimary!75] (axis cs:102.8936,0.0706) rectangle (axis cs:126.8936,6.6384);
\fill[PlotSecondary!35] (axis cs:102.8936,6.7090) rectangle (axis cs:126.8936,8.4746);
\fill[PlotPrimary!75] (axis cs:128.4255,0.0706) rectangle (axis cs:152.4255,4.4492);
\fill[PlotSecondary!35] (axis cs:128.4255,4.5198) rectangle (axis cs:152.4255,7.8390);
\fill[PlotPrimary!75] (axis cs:153.9574,0.0706) rectangle (axis cs:177.9574,2.9661);
\fill[PlotSecondary!35] (axis cs:153.9574,3.0367) rectangle (axis cs:177.9574,4.9435);
\fill[PlotPrimary!75] (axis cs:179.4894,0.0706) rectangle (axis cs:203.4894,2.3305);
\fill[PlotSecondary!35] (axis cs:179.4894,4.4492) rectangle (axis cs:203.4894,4.7316);
\fill[PlotSecondary!35] (axis cs:179.4894,2.4011) rectangle (axis cs:203.4894,4.2373);
\fill[PlotPrimary!75] (axis cs:205.0213,0.0706) rectangle (axis cs:229.0213,1.8362);
\fill[PlotSecondary!35] (axis cs:205.0213,1.9068) rectangle (axis cs:229.0213,3.8842);
\fill[PlotPrimary!75] (axis cs:230.5532,0.0706) rectangle (axis cs:254.5532,1.4124);
\fill[PlotSecondary!35] (axis cs:230.5532,1.4831) rectangle (axis cs:254.5532,2.8249);
\fill[PlotPrimary!75] (axis cs:256.0851,0.0706) rectangle (axis cs:280.0851,1.0593);
\fill[PlotSecondary!35] (axis cs:256.0851,1.1299) rectangle (axis cs:280.0851,2.5424);
\fill[PlotPrimary!75] (axis cs:281.6170,0.0706) rectangle (axis cs:305.6170,1.2712);
\fill[PlotSecondary!35] (axis cs:281.6170,1.3418) rectangle (axis cs:305.6170,3.6723);
\fill[PlotPrimary!75] (axis cs:307.1489,0.0706) rectangle (axis cs:331.1489,0.7768);
\fill[PlotSecondary!35] (axis cs:307.1489,0.8475) rectangle (axis cs:331.1489,1.8362);
\fill[PlotPrimary!75] (axis cs:332.6809,0.0706) rectangle (axis cs:356.6809,0.6356);
\fill[PlotSecondary!35] (axis cs:332.6809,0.7062) rectangle (axis cs:356.6809,1.5537);
\fill[PlotPrimary!75] (axis cs:358.2128,0.0706) rectangle (axis cs:382.2128,0.4944);
\fill[PlotSecondary!35] (axis cs:358.2128,0.5650) rectangle (axis cs:382.2128,1.4124);
\fill[PlotPrimary!75] (axis cs:383.7447,0.0706) rectangle (axis cs:407.7447,0.4237);
\fill[PlotSecondary!35] (axis cs:383.7447,0.4944) rectangle (axis cs:407.7447,1.4124);
\fill[PlotPrimary!75] (axis cs:409.2766,0.0706) rectangle (axis cs:433.2766,0.2825);
\fill[PlotSecondary!35] (axis cs:409.2766,0.3531) rectangle (axis cs:433.2766,0.9887);
\fill[PlotPrimary!75] (axis cs:434.8085,0.0706) rectangle (axis cs:458.8085,0.2825);
\fill[PlotSecondary!35] (axis cs:434.8085,0.3531) rectangle (axis cs:458.8085,0.9181);
\fill[PlotPrimary!75] (axis cs:460.3404,0.0706) rectangle (axis cs:484.3404,0.2119);
\fill[PlotSecondary!35] (axis cs:460.3404,0.2825) rectangle (axis cs:484.3404,0.7062);
\fill[PlotPrimary!75] (axis cs:485.8723,0.0706) rectangle (axis cs:509.8723,0.1412);
\fill[PlotSecondary!35] (axis cs:485.8723,0.2119) rectangle (axis cs:509.8723,1.2712);
\fill[PlotPrimary!75] (axis cs:511.4043,0.0706) rectangle (axis cs:535.4043,0.1412);
\fill[PlotSecondary!35] (axis cs:511.4043,0.2119) rectangle (axis cs:535.4043,0.6356);
\fill[PlotSecondary!35] (axis cs:536.9362,0.1412) rectangle (axis cs:560.9362,0.4944);
\fill[PlotSecondary!35] (axis cs:562.4681,0.1412) rectangle (axis cs:586.4681,0.4237);
\fill[PlotSecondary!35] (axis cs:588.0000,0.1412) rectangle (axis cs:612.0000,0.3531);
\fill[PlotSecondary!35] (axis cs:613.5319,0.1412) rectangle (axis cs:637.5319,0.4237);
\fill[PlotSecondary!35] (axis cs:639.0638,0.0706) rectangle (axis cs:663.0638,0.2825);
\fill[PlotSecondary!35] (axis cs:664.5957,0.0706) rectangle (axis cs:688.5957,0.2119);
\fill[PlotSecondary!35] (axis cs:690.1277,0.0706) rectangle (axis cs:714.1277,0.2119);
\fill[PlotSecondary!35] (axis cs:715.6596,0.0706) rectangle (axis cs:739.6596,0.2119);
\fill[PlotSecondary!35] (axis cs:741.1915,0.0706) rectangle (axis cs:765.1915,0.4237);
\fill[PlotSecondary!35] (axis cs:766.7234,0.0706) rectangle (axis cs:790.7234,0.1412);
\fill[PlotSecondary!35] (axis cs:792.2553,0.0706) rectangle (axis cs:816.2553,0.1412);
\fill[PlotSecondary!35] (axis cs:817.7872,0.0706) rectangle (axis cs:841.7872,0.1412);
\addlegendimage{area legend,fill=PlotSecondary!35}\addlegendentry{Moderate PA}\addlegendimage{area legend,fill=PlotPrimary!75}\addlegendentry{Vigorous PA}\end{axis}\end{tikzpicture}\hfill\begin{tikzpicture}\begin{axis}[width=7cm,height=5.4cm,xmin=0,xmax=900,ymin=.4,ymax=1.2,xtick={0,150,300,450,600,750,900},ytick={.4,.6,.8,1,1.2},xlabel={Physical Activity (min/week)},ylabel={HR (95\%CI)},title={Non-CVD mortality},axis lines=left,title style={font=\small\bfseries},legend style={font=\scriptsize},legend pos=south west]
\addplot[PlotSecondary,line width=1pt,dashed] coordinates {(0.000000,1.000000) (19.251337,0.970667) (25.668449,0.971733) (32.085561,0.963200) (38.502674,0.953600) (44.919786,0.945067) (51.336898,0.937600) (57.754011,0.929067) (64.171123,0.922667) (70.588235,0.915200) (77.005348,0.908800) (83.422460,0.902400) (89.839572,0.896000) (96.256684,0.891733) (102.673797,0.884267) (109.090909,0.874667) (115.508021,0.870400) (121.925134,0.866133) (128.342246,0.855467) (134.759358,0.852267) (141.176471,0.848000) (147.593583,0.844800) (154.010695,0.840533) (160.427807,0.838400) (166.844920,0.834133) (173.262032,0.829867) (179.679144,0.827733) (186.096257,0.821333) (192.513369,0.818133) (198.930481,0.817067) (205.347594,0.813867) (211.764706,0.809600) (218.181818,0.808533) (224.598930,0.804267) (231.016043,0.803200) (237.433155,0.801067) (243.850267,0.798933) (250.267380,0.795733) (256.684492,0.794667) (263.101604,0.792533) (269.518717,0.791467) (275.935829,0.789333) (282.352941,0.788267) (288.770053,0.784000) (295.187166,0.784000) (301.604278,0.782933) (308.021390,0.781867) (314.438503,0.778667) (320.855615,0.778667) (327.272727,0.778667) (333.689840,0.776533) (340.106952,0.776533) (346.524064,0.774400) (352.941176,0.773333) (359.358289,0.772267) (365.775401,0.770133) (372.192513,0.769067) (378.609626,0.769067) (385.026738,0.768000) (391.443850,0.765867) (397.860963,0.763733) (404.278075,0.763733) (410.695187,0.763733) (417.112299,0.762667) (423.529412,0.761600) (429.946524,0.760533) (436.363636,0.759467) (442.780749,0.758400) (449.197861,0.757333) (455.614973,0.756267) (462.032086,0.756267) (468.449198,0.754133) (474.866310,0.752000) (481.283422,0.752000) (487.700535,0.750933) (494.117647,0.750933) (500.534759,0.747733) (506.951872,0.745600) (513.368984,0.745600) (519.786096,0.745600) (526.203209,0.744533) (532.620321,0.744533) (539.037433,0.742400) (545.454545,0.742400) (551.871658,0.740267) (558.288770,0.739200) (564.705882,0.738133) (571.122995,0.738133) (577.540107,0.737067) (583.957219,0.734933) (590.374332,0.732800) (596.791444,0.733867) (603.208556,0.732800) (609.625668,0.731733) (616.042781,0.728533) (622.459893,0.728533) (628.877005,0.727467) (635.294118,0.727467) (641.711230,0.726400) (648.128342,0.725333) (654.545455,0.724267) (660.962567,0.723200) (667.379679,0.722133) (673.796791,0.721067) (680.213904,0.718933) (686.631016,0.720000) (693.048128,0.718933) (699.465241,0.716800) (705.882353,0.714667) (712.299465,0.715733) (718.716578,0.713600) (725.133690,0.711467) (731.550802,0.710400) (737.967914,0.711467) (744.385027,0.710400) (750.802139,0.710400) (757.219251,0.708267) (763.636364,0.708267) (770.053476,0.707200) (776.470588,0.705067) (782.887701,0.705067) (789.304813,0.704000) (795.721925,0.702933) (802.139037,0.701867) (808.556150,0.700800) (814.973262,0.697600) (821.390374,0.697600) (827.807487,0.698667) (834.224599,0.696533) (840.641711,0.695467) (847.058824,0.694400) (853.475936,0.693333) (859.893048,0.693333) (866.310160,0.692267) (872.727273,0.691200) (879.144385,0.690133) (885.561497,0.690133) (891.978610,0.688000) (898.395722,0.688000)};\addlegendentry{Moderate PA, $P$ for non-linear $<0.001$}
\addplot[PlotPrimary,line width=1pt] coordinates {(0.000000,1.000000) (19.251337,0.970667) (25.668449,0.960000) (32.085561,0.950400) (38.502674,0.940800) (44.919786,0.932267) (51.336898,0.923733) (57.754011,0.916267) (64.171123,0.908800) (70.588235,0.902400) (77.005348,0.897067) (83.422460,0.891733) (89.839572,0.886400) (96.256684,0.881067) (102.673797,0.877867) (109.090909,0.873600) (115.508021,0.870400) (121.925134,0.866133) (128.342246,0.864000) (134.759358,0.860800) (141.176471,0.858667) (147.593583,0.856533) (154.010695,0.854400) (160.427807,0.854400) (166.844920,0.852267) (173.262032,0.852267) (179.679144,0.851200) (186.096257,0.850133) (192.513369,0.848000) (198.930481,0.848000) (205.347594,0.846933) (211.764706,0.845867) (218.181818,0.845867) (224.598930,0.844800) (231.016043,0.844800) (237.433155,0.844800) (243.850267,0.843733) (250.267380,0.843733) (256.684492,0.843733) (263.101604,0.842667) (269.518717,0.842667) (275.935829,0.841600) (282.352941,0.841600) (288.770053,0.840533) (295.187166,0.839467) (301.604278,0.839467) (308.021390,0.839467) (314.438503,0.838400) (320.855615,0.837333) (327.272727,0.837333) (333.689840,0.837333) (340.106952,0.836267) (346.524064,0.836267) (352.941176,0.836267) (359.358289,0.835200) (365.775401,0.835200) (372.192513,0.834133) (378.609626,0.834133) (385.026738,0.833067) (391.443850,0.832000) (397.860963,0.830933) (404.278075,0.830933) (410.695187,0.830933) (417.112299,0.829867) (423.529412,0.828800) (429.946524,0.828800) (436.363636,0.828800) (442.780749,0.828800) (449.197861,0.828800) (455.614973,0.827733) (462.032086,0.826667) (468.449198,0.826667) (474.866310,0.826667) (481.283422,0.825600) (487.700535,0.824533) (494.117647,0.824533) (500.534759,0.822400) (506.951872,0.822400) (513.368984,0.822400) (519.786096,0.822400) (526.203209,0.822400) (532.620321,0.821333) (539.037433,0.820267) (545.454545,0.820267) (551.871658,0.820267) (558.288770,0.819200) (564.705882,0.818133) (571.122995,0.818133) (577.540107,0.818133) (583.957219,0.818133) (590.374332,0.818133) (596.791444,0.817067) (603.208556,0.816000) (609.625668,0.816000) (616.042781,0.814933) (622.459893,0.813867) (628.877005,0.813867) (635.294118,0.813867) (641.711230,0.813867) (648.128342,0.812800) (654.545455,0.811733) (660.962567,0.811733) (667.379679,0.811733) (673.796791,0.810667) (680.213904,0.810667) (686.631016,0.809600) (693.048128,0.809600) (699.465241,0.809600) (705.882353,0.809600) (712.299465,0.807467) (718.716578,0.808533) (725.133690,0.805333) (731.550802,0.805333) (737.967914,0.805333) (744.385027,0.805333) (750.802139,0.804267) (757.219251,0.804267) (763.636364,0.803200) (770.053476,0.803200) (776.470588,0.803200) (782.887701,0.803200) (789.304813,0.802133) (795.721925,0.802133) (802.139037,0.801067) (808.556150,0.801067) (814.973262,0.800000) (821.390374,0.800000) (827.807487,0.798933) (834.224599,0.798933) (840.641711,0.797867) (847.058824,0.796800) (853.475936,0.796800) (859.893048,0.796800) (866.310160,0.796800) (872.727273,0.795733) (879.144385,0.794667) (885.561497,0.794667) (891.978610,0.794667) (898.395722,0.793600)};\addlegendentry{Vigorous PA, $P$ for non-linear $<0.001$}
\end{axis}\end{tikzpicture}}
限制性立方样条模型显示，VPA与全因死亡、心血管疾病和非心血管疾病的死亡风险显著降低相关，但当VPA达到更长时间时，相关性趋于平稳（图中实线；原文图中为黑实线）。限制性立方样条模型显示，MPA与全因死亡、心血管疾病和非心血管疾病的死亡风险显著降低相关，随着时长的增加，相关性在增强（图中虚线；原文图中为灰实线）。
\begin{annotation}
以“不运动”（0 min/wk）为参照读图：全因死亡中，VPA的HR在约150 min/wk以内下降较快（约0.85），之后几乎持平（900 min/wk时约0.80），即超过指南推荐量后增加剧烈运动的额外获益很小；MPA的HR在整个观察范围内持续下降（150 min/wk约0.84，900 min/wk约0.65）。CVD死亡中，VPA曲线在150--300 min/wk处最低（约0.84），之后略有回升（900 min/wk约0.89），但该段观察人数少、区间宽，不足以断定剧烈运动过多会减弱保护作用。两条曲线的非线性检验均 $P<0.001$，说明剂量反应关系不是对数线性的。时长很大处观察人数少（见左下直方图），曲线尾部的精度较低。
\end{annotation}
\subsection{英文表达}
\sourcepages{6}{53-55}
\textbf{方法}
\begin{quote}
METHODS: A total of 116 221 adults from 2 large prospective US cohorts (Nurses' Health Study and Health Professionals Follow-up Study, 1988--2018) were analyzed. Detailed self-reported leisure-time physical activity was assessed with a validated questionnaire, repeated up to 15 times during the follow-up. Cox regression was used to estimate the hazard ratio and 95\% CI of the association between long-term leisure-time physical activity intensity and all-cause and cause-specific mortality.
\end{quote}
\textbf{结果}
\begin{quote}
During 30 years of follow-up, we identified 47 596 deaths.

Compared with those with no long-term leisure-time VPA, participants who met the long-term leisure-time physical activity guidelines (75--149 min/wk of VPA) had 19\% lower all-cause mortality (HR, 0.81 [95\% CI, 0.76--0.87]), 31\% lower CVD mortality (HR, 0.69 [95\% CI, 0.60--0.78]), and 15\% lower non-CVD mortality (HR, 0.85 [95\% CI, 0.79--0.92]; Table 3).

In the restricted cubic spline models, long-term leisure-time MPA was associated with substantially lower risk of all-cause, CVD, and non-CVD mortality, and the inverse association for long-term leisure-time MPA continued to strengthen over the entire range of activity assessed in the cohort (Figure 1).
\end{quote}
\subsection{注意事项}
\sourcepages{6}{56}
效应值的计算和解释；前提假定（比例风险假定）；人年的计算。
\section{思考与练习}
\sourcepages{6}{57-59}
\textbf{1.} 比较两种治疗疗法对延长肺癌患者存活时间的效果。观察了40例接受标准治疗和试验疗法的病人，同时记录了患者的年龄、病程、病理类型等信息。具体数据见survivalhomework1。

变量信息：time（生存时间，天），status（结局死亡1，删失0），treat（标准疗法0，试验疗法1），age（年龄），coursetime（病程），type（病理类型，1鳞状、2小型、3腺状、4大型）。
\begin{enumerate}
\item 绘制并比较两个治疗组的生存曲线。
\item 在控制了年龄、病程和病理类型之后，比较两个治疗组的疗效。
\item 哪些因素影响患者的存活时间？
\item 如何评价所使用的模型是否合适？所构建的模型是较优的？
\item 如何解释结果？不同治疗方法的疗效比较结果为差别没有统计学意义时，如何理解？
\end{enumerate}
\sourcepages{6}{60-61}
\textbf{2. 综合分析题}

\textit{Mortality risk prediction of high-sensitivity C-reactive protein in suspected acute coronary syndrome: A cohort study}。

Amit Kaura, Adam Hartley, Vasileios Panoulas, Ben Glampson, Anoop S V Shah, Jim Davies, Abdulrahim Mulla, Kerrie Woods, Joe Omigie, Anoop D Shah, Mark R Thursz, Paul Elliott, Harry Hemingway, Bryan Williams, Folkert W Asselbergs, Michael O'Sullivan, Graham M Lord, Adam Trickey, Jonathan A C Sterne, Dorian O Haskard, Narbeh Melikian, Darrel P Francis, Wolfgang Koenig, Ajay M Shah, Rajesh Kharbanda, Divaka Perera, Riyaz S Patel, Keith M Channon, Jamil Mayet, Ramzi Khamis.

PLoS Med. 2022 Feb 22;19(2):e1003911. doi:10.1371/journal.pmed.1003911. eCollection 2022 Feb. PMID:35192610. PMCID:PMC8863282。
\begin{enumerate}
\item 请凝练本篇文章要解决的科学问题。
\item 原始数据结构主要包括了哪些变量？变量类型如何？
\item 这篇文章的流行病学设计是什么？基于此设计，请简述如何用统计模型（Cox回归模型）来解决科学问题。
\end{enumerate}
\par\Needspace{8\baselineskip}\noindent\textbf{补充练习（诊断题）}
\question{某研究比较心脏移植对生存的影响：以进入移植等待名单为时间起点，把“随访中接受过移植”作为基线0/1变量放入Cox模型，得到HR $=0.40$。这一分析有什么问题？应如何改正？}
\answer{接受移植的人必须先在等待期内存活到移植那一天，这段时间被错误地归入“移植组”，形成不朽时间偏倚，使移植组显得风险更低。应把移植作为时依协变量：移植前的随访时间计入未移植状态，移植后计入移植状态，用 $(\text{start},\text{stop}]$ 计数过程数据拟合Cox模型；也可用界标分析（landmark analysis）。本章表8中的relapse存在同样的问题。}
\section*{小结}
\sourcepages{6}{62-63}
生存分析的目的、资料特点、基本概念；生存率的估计方法（寿命表法、Kaplan--Meier）；生存曲线比较的方法（log-rank法）；Cox回归分析；生存分析方法的正确应用。
\section{附录：软件实现}
\sourcepages{6}{64-65}
\Needspace{9\baselineskip}
\subsection{R语言的实现}
\par\noindent
\begin{coursecode}{R}
library(survival)
fit.cox <- coxph(Surv(time, status) ~ x1+x2+x3+x4, data=mydata)
summary(fit.cox)
basehaz(fit.cox, centered=FALSE)   # 协变量取0时的累积基线风险
phtest <- cox.zph(fit.cox)
\end{coursecode}
\par
\begin{annotation}函数名为 \texttt{Surv}（首字母大写，\texttt{surv} 不存在）。\texttt{status} 中1为事件、0为删失（也可用TRUE/FALSE）；分类变量用 \texttt{factor()}，以 \texttt{relevel()} 指定参照组。\texttt{coxph} 默认用Efron法处理并列事件，可用 \texttt{ties="breslow"} 与SAS/SPSS对齐。\texttt{cox.zph} 用Schoenfeld残差检验比例风险假设；时依协变量用 \texttt{Surv(start, stop, status)} 形式的计数过程数据。以第4章例3验证：\texttt{coxph(Surv(t, e) \textasciitilde{} size)} 得 $\widehat\beta=1.275$（HR 3.58），\texttt{survdiff} 得log-rank $\chi^2=6.77$，均与正文一致。\end{annotation}
\subsection{SPSS软件实现}
\sourcepages{6}{66}
Analyze $\to$ Survival：Life table、Kaplan--Meier、Cox regression、Cox w/Time-Dep Cov。
\Needspace{8\baselineskip}
\subsection{SAS软件实现}
\sourcepages{6}{67}
\par\noindent
\begin{coursecode}{SAS}
PROC PHREG;
MODEL t*status(0)=age grade size relapse/
  SELECTION=STEPWISE SLE=0.05 SLS=0.05;
RUN;
\end{coursecode}
\par
\Needspace{8\baselineskip}
\subsection{STATA软件实现}
\sourcepages{6}{68}
\par\noindent
\begin{coursecode}{Stata}
stset TIME, failure(STATUS==1) // 定义时间和结局变量，STATUS为1表示事件
stcox x1 x2 x3 x4                // 定义模型中的自变量（默认输出HR）
stcurve, survival                // 绘制生存曲线
estat phtest, detail             // 基于Schoenfeld残差检验PH假定
\end{coursecode}
\par

\begin{fuxi}[复习题]
\begin{enumerate}
\item 说明Cox模型称为半参数模型的原因。\\
\textit{答：协变量部分$\exp(\bm x^\top\bm\beta)$有参数形式，基线风险$h_0(t)$不作任何分布假设。}
\item 肿瘤分级按1、2、3赋值，$\widehat\beta=1.572$，解释HR的含义。\\
\textit{答：其他因素不变时，分级每提高1级，死亡风险为原来的$e^{1.572}=4.82$倍；该解释假定相邻两级之间的HR相同。}
\item 说明部分似然中不出现$h_0(t)$的原因。\\
\textit{答：在每个事件时刻，风险集中恰为该个体发生事件的条件概率中，$h_0(t)$在分子与分母中相消。}
\item 列出检查比例风险假定的方法。\\
\textit{答：观察KM曲线是否交叉，log--log图是否平行；Schoenfeld残差检验（\texttt{cox.zph}）；在模型中加入协变量与时间的交互项。}
\item 列出比例风险假定不成立时的处理方法。\\
\textit{答：分层Cox模型，分时段拟合，加入协变量与时间函数的交互项（时变系数）。}
\item 患者在随访中复发，说明能否将是否复发作为基线协变量纳入Cox模型。\\
\textit{答：不能，否则会产生不朽时间偏倚；应记录复发时间，按时依协变量处理。}
\end{enumerate}
\end{fuxi}
